\documentclass[12pt]{article}

\usepackage{amssymb}
\usepackage{amsfonts}
\usepackage{amsmath}
\usepackage[nohead]{geometry}
\usepackage[onehalfspacing]{setspace}
\usepackage[bottom]{footmisc}
\usepackage{indentfirst}
\usepackage{endnotes}
\usepackage{threeparttable}
\usepackage{multirow}
\usepackage{dsfont}
\usepackage{graphicx}
\usepackage[hang]{caption}
\usepackage{lscape}
\usepackage{natbib}
\usepackage{enumerate}
\usepackage{color} \definecolor{orange}{RGB}{255,127,0} \definecolor{brown}{RGB}{150,75,0}
\usepackage{subcaption}
\usepackage{appendix}
\usepackage{calc}
\usepackage{pdflscape}
\usepackage{afterpage}
\usepackage{array}
\usepackage{tabularx}
\graphicspath{ {./Figures/} }
\usepackage{colortbl}
\usepackage{rotating}
\usepackage{empheq}
\usepackage{soul} %crosses out text
\usepackage{titlesec}
\usepackage{xfrac} % for fractions in text 1/4
\usepackage{accents} % to do underbars
\usepackage{afterpage} % to prevent empty spaces before landscape
\usepackage{authblk} % to customize title page
\usepackage{xr-hyper}% to add external references from the Appendix
\usepackage[normalem]{ulem}
\makeatletter
\newcommand*{\addFileDependency}[1]{% argument=file name and extension
  \typeout{(#1)}
  \@addtofilelist{#1}
  \IfFileExists{#1}{}{\typeout{No file #1.}}
}
\makeatother
\newcommand*{\myexternaldocument}[1]{%
    \externaldocument{#1}%
    \addFileDependency{#1.tex}%
    \addFileDependency{#1.aux}%
}
\myexternaldocument{Appendix}

\usepackage[colorlinks=true,urlcolor=rblue,citecolor=rblue,linkcolor=rblue,bookmarks=true]{hyperref}
\usepackage{dcolumn} % to align periods in tables
\newcolumntype{d}[1]{D{.}{.}{#1}}

\usepackage{url}
\def\UrlBreaks{\do\/\do-} % to brake urls across lines and avoid overflow

% the following block is to prevent weird hyphenation of words across lines
\pretolerance=10000
\tolerance=2000
\emergencystretch=3em
\hyphenpenalty=5000
\exhyphenpenalty=5000
\sloppy

\newtheorem{theorem}{Theorem}
\newtheorem{acknowledgement}{Acknowledgement}
\newtheorem{algorithm}[theorem]{Algorithm}
\newtheorem{axiom}[theorem]{Axiom}
\newtheorem{case}[theorem]{Case}
\newtheorem{claim}[theorem]{Claim}
\newtheorem{conclusion}[theorem]{Conclusion}
\newtheorem{condition}[theorem]{Condition}
\newtheorem{conjecture}[theorem]{Conjecture}
\newtheorem{corollary}[theorem]{Corollary}
\newtheorem{criterion}[theorem]{Criterion}
\newtheorem{definition}[theorem]{Definition}
\newtheorem{example}[theorem]{Example}
\newtheorem{exercise}[theorem]{Exercise}
\newtheorem{lemma}[theorem]{Lemma}
\newtheorem{notation}[theorem]{Notation}
\newtheorem{problem}[theorem]{Problem}
\newtheorem{proposition}{Proposition}
\newtheorem{remark}[theorem]{Remark}
\newtheorem{solution}[theorem]{Solution}
\newtheorem{summary}[theorem]{Summary}
\makeatletter
\def\@biblabel#1{\hspace*{-\labelsep}}
\makeatother
\geometry{left=1in,right=1in,top=1.00in,bottom=1.0in}
\bibpunct{(}{)}{;}{a}{,}{,}

\definecolor{rblue}{cmyk}{0.71,0.53,0,0.12}
\setlength{\captionmargin}{50pt}
\setcounter{MaxMatrixCols}{30}
\newcommand{\E}{\mathbb{E}}
\newcommand{\eps}{\varepsilon}
\newcommand{\h}{\hat}
\newcommand{\bp}{\bar{\pi}}
\newcommand{\bP}{\bar{\Pi}}
\newcommand{\bR}{\bar{R}}
\newcommand{\bb}{\bar{b}}
\newcommand{\btau}{\bar{\tau}}
\newcommand{\vphi}{\varphi}
\newcommand{\tP}{\tilde{P}}
\newcommand{\tp}{\tilde{p}}
\newcommand{\hA}{\hat{A}}
\newcommand{\hg}{\hat{g}}
\newcommand{\hc}{\hat{c}}
\newcommand{\hy}{\hat{y}}
\newcommand{\hY}{\hat{Y}}
\newcommand{\hK}{\hat{K}}
\newcommand{\hp}{\hat{\pi}}
\newcommand{\hpsi}{\hat{\psi}}
\newcommand{\hR}{\hat{R}}
\newcommand{\htau}{\hat{\tau}}
\newcommand{\hP}{\hat{\Pi}}
\newcommand{\hmp}{\hat{\left(\frac{m}{p}\right)}}
\newcommand{\hb}{\hat{b}}
\newcommand{\hd}{\hat{d}}
\newcommand{\hthe}{\hat{\theta}}
\newcommand{\hnu}{\hat{\nu}}
\newcommand{\hv}{\hat{v}}
\newcommand{\tpi}{\tilde{\pi}}
\newcommand{\F}{\mathfrak{F}}
\newcommand{\Z}{\mathfrak{Z}}
\newcommand{\Hh}{\mathfrak{H}}
\newcommand{\w}{\omega}
\newcommand{\Rho}{\textrm{P}}
\newcommand{\Mu}{\textrm{M}}
\newcommand{\e}{\mathrm{e}}
\newcommand{\ex}[2]{\ensuremath{\e^{- {#1} (x - {#2})}}}
\newcommand{\ey}[2]{\ensuremath{\e^{- {#1} (y - {#2})}}}
\newcommand{\exy}[4]{\ensuremath{\e^{- {#1} (x - {#2}) - {#3} (y - {#4})}}}
\newcommand{\al}{\alpha}
\newcommand{\be}{\beta}
\newcommand{\de}{\delta}
\newcommand{\bbeta}{\boldsymbol{\beta}}
\newcommand{\bups}{\boldsymbol{\upsilon}}
\newcommand{\btheta}{\boldsymbol{\theta}}
\newcommand{\bXttm}{\mathbf{X}_{t|t-1}}
\newcommand{\bXtmtm}{\mathbf{X}_{t-1|t-1}}
\newcommand{\bXtt}{\mathbf{X}_{t|t}}
\newcommand{\bYttm}{\mathbf{Y}_{t|t-1}}
\newcommand{\bYtt}{\mathbf{Y}_{t|t}}
\newcommand{\bOmegattm}{\mathbf{\Omega}_{t|t-1}}
\newcommand{\bOmegatmtm}{\mathbf{\Omega}_{t-1|t-1}}
\newcommand{\bOmegatt}{\mathbf{\Omega}_{t|t}}
\newcommand{\bSigmat}{\mathbf{\Sigma}_t}
\newcommand{\BSigmatt}{\mathbf{\Sigma}_{t|t}}
\newcommand{\balpha}{\boldsymbol{\alpha}}
\newcommand{\bgamma}{\boldsymbol{\gamma}}
\newcommand{\brho}{\boldsymbol{\rho}}
\newcommand{\bT}{\mathbf{T}}
\newcommand{\Prob}{\mathbb{P}}
\newcommand{\head}[1]{\textnormal{\textbf{#1}}}
\newcommand{\bof}[1]{\mathbf{#1}}
\newcommand{\bos}[1]{\boldsymbol{#1}}
\newcommand{\dy}[1]{\Delta y_{#1}}
\newcommand{\du}[1]{\Delta u_{#1}}
\newcolumntype{C}{>{\centering\arraybackslash} m{3cm} }
\renewcommand{\appendixpagename}{Appendix}
\DeclareMathOperator{\cov}{cov}
\DeclareMathOperator{\var}{var}
\DeclareMathOperator{\corr}{corr}
\DeclareMathOperator{\vecc}{vec}
\allowdisplaybreaks
\newcommand{\ubar}[1]{\underaccent{\bar}{#1}} % to do underbars

\begin{document}

\title{Fiscal Consolidations in Commodity-Exporting Countries: A Small Open Economy DSGE Perspective\thanks{The views expressed in this paper are solely the responsibility of the authors and should not be interpreted as reflecting the views of the Board of Governors of the Federal Reserve System or the Central Bank of Chile.}}

\author{
Manuel González-Astudillo\thanks{Board of Governors of the Federal Reserve System and Escuela Superior Polit\'ecnica del Litoral. E-mail address: \url{manuel.p.gonzalez-astudillo@frb.gov}.}
$\qquad$
Juan Guerra-Salas\thanks{Corresponding author. Central Bank of Chile. E-mail address: \url{jguerra@bcentral.cl}.}
$\qquad$
Avi Lipton\thanks{Harvard University. E-mail address: \url{alipton@g.harvard.edu}.}
}

\date{\today}
\maketitle

\begin{abstract}
\singlespacing
%  words
Fiscal consolidation in commodity‑exporting economies is subject to substantial uncertainty because fiscal revenues depend heavily on volatile international commodity prices. This paper quantifies how deviations in commodity revenues affect consolidation outcomes using a dynamic stochastic general equilibrium model for a small open economy with a commodity sector, distortionary taxation, productive public spending, heterogeneous households, and external borrowing subject to an endogenous sovereign risk premium. The model is estimated using Ecuadorian data for 2004–2019 and then used to simulate Ecuador’s 2020–2025 IMF‑supported fiscal consolidation under perfect foresight, treating the program’s fiscal instruments as fully anticipated policy paths. We find that historically plausible deviations in oil revenues generate pronounced asymmetries in fiscal and macroeconomic dynamics. Higher‑than‑expected revenues accelerate deficit reduction, lower sovereign spreads, and stimulate economic activity, potentially weakening incentives to sustain adjustment, while revenue shortfalls raise borrowing costs, depress investment, and intensify the required fiscal effort. These results highlight the importance of explicitly accounting for commodity‑revenue uncertainty when evaluating consolidation programs. While the analysis is calibrated to Ecuador, the framework and qualitative findings are broadly applicable to commodity‑exporting economies implementing fiscal adjustment under external financing constraints.


\end{abstract}

\textit{Keywords}: Fiscal consolidation, commodity exporter, small open economy, DSGE model


\textit{JEL classification}: E62, F34, F47, Q43

%%%%%%%%%%%%%%%%%%%%%%%%%%%%%%%%%%%%%%%%%%%%%%%%%%%%%%

\thispagestyle{empty}
\setcounter{page}{0}
\clearpage
\setcounter{page}{1}

%%%%%%%%%%%%%%%%%%%%%%%%%%%%%%%%%%%%%%%%%%%%%%%%%%%%%%

\section{Introduction}

When the boom in commodity prices ended around 2014, several commodity-exporting countries faced deteriorated fiscal sustainability prospects and higher risks of debt distress \citep[see][]{richaud2019fiscal}. The COVID-19 pandemic exacerbated these problems, with sovereign credit spreads such as the Emerging Market Bond Index Global (EMBIG) increasing by about 300 bps. Moreover, spreads increased by more than 100 bps after the turbulence caused by the global monetary tightening cycle that followed the post-pandemic rise in inflation, and the Russian invasion of Ukraine. These tighter financial conditions affected the fiscal situation of many countries to the point that they needed to ask international organizations for extraordinary financing in order to face these challenges. Indeed, the IMF lists 90 countries currently receiving or having received about USD 170 billion in financial assistance and debt service relief from March 2020 to March 2022.\footnote{See \url{https://www.imf.org/en/Topics/imf-and-covid19/COVID-Lending-Tracker}.} This type of funding is typically tied to the implementation of fiscal consolidation programs aimed at putting public finances on a sustainable trajectory. 

In commodity-exporting countries, fiscal consolidation plans are particularly challenging to implement due to the high dependence of fiscal revenue on volatile international commodity prices. This volatility introduces significant uncertainty into revenue forecasts, complicating the design and implementation of credible consolidation strategies. To study these risks, we develop a small open economy dynamic stochastic general equilibrium (SOE DSGE) model with a commodity sector, distortionary taxes, productive public spending, heterogeneous households, and external borrowing subject to an endogenous sovereign risk premium. This model provides a structural framework to analyze how deviations in commodity revenues affect fiscal balances and macroeconomic outcomes during consolidation episodes.

We estimate the model using Ecuadorian quarterly data for 2004:Q1\textendash 2019:Q4 and then evaluate a concrete IMF‑supported fiscal consolidation program: Ecuador’s 2020–2025 Extended Fund Facility adjustment package. The policy experiment is implemented under perfect foresight by solving the nonlinear model, treating the program’s fiscal instruments\textemdash tax rates, government consumption, public investment, and transfers\textemdash as fully anticipated deterministic paths. This design incorporates anticipation effects and preserves the nonlinear dynamics of the program simulation, while maintaining a clear link between the estimated model and the policy counterfactual.

Our main finding is that historically plausible deviations in commodity revenues generate pronounced asymmetries in consolidation outcomes. When oil revenue exceeds program projections, the fiscal balance improves faster, the risk premium falls, and economic activity strengthens\textemdash potentially weakening incentives to sustain fiscal effort in practice. When oil revenue falls short, the fiscal balance deteriorates relative to baseline, spreads rise, and private investment weakens, implying that meeting program targets would require additional adjustment. These asymmetric dynamics operate through the interaction of revenue realizations with external debt and sovereign spreads, which affect domestic financing conditions during the consolidation process. The results underscore the value of accounting for revenue uncertainty when assessing consolidation feasibility.

Ecuador provides a particularly informative environment because the 2020–2025 program specifies detailed paths for multiple fiscal instruments together with oil‑revenue projections, allowing the program to be mapped into explicit policy paths in the model. More broadly, the paper’s qualitative mechanism applies to other commodity exporters that consolidate under external borrowing constraints and volatile sovereign spreads: consolidation performance is highly contingent on commodity‑revenue realizations and the financing conditions they induce.

The paper contributes to the literature in three ways. First, it evaluates a multi-instrument IMF-style fiscal consolidation package using a fully specified structural SOE DSGE model estimated on pre-program data. This approach differs from standard financial programming or macroframework projection tools typically used in policy institutions and allows us to trace nonlinear interactions among commodity revenue realizations, external debt dynamics, and sovereign risk premia under perfect foresight. Second, it highlights a macro‑financial channel through which commodity‑revenue deviations shift debt and sovereign spreads, shaping output and investment during consolidation. Third, it quantifies consolidation risk in a realistic policy setting by simulating the program under alternative revenue paths calibrated to historical volatility, showing how the same fiscal measures can succeed, overshoot, or fail depending on commodity revenues.

The remainder of the paper is organized as follows: Section \ref{sec:literature} reviews the related literature; Section \ref{sec:model} presents the model; Section \ref{sec:parametrization} discusses calibration and Bayesian estimation; Section \ref{sec:performance} evaluates model fit and impulse responses; Section \ref{sec:fiscalpolicy} analyzes fiscal consolidations under perfect foresight and alternative oil‑revenue paths; and Section \ref{sec:conclusion} concludes.

%%%%%%%%%%%%%%%%%%%%%%%%%%%%%

\section{Related Literature}\label{sec:literature}

%%%%%%%%%%%%%%%%%%%%%%%%%%%%%

Evaluating fiscal consolidations in economies that rely heavily on commodity revenues is inherently risky. The central challenge is that consolidation outcomes often depend less on the design of the adjustment itself than on the trajectory of volatile commodity-related revenue. In such settings, the success or failure of consolidation can be confounded by commodity price shocks, making ex ante assessments highly uncertain.

A large body of work documents how oil and other commodity revenues dominate macroeconomic dynamics in resource exporters. \citet{GarciaAlban2021} show that in Ecuador, oil revenue shocks have historically been the most important driver of output fluctuations, far outweighing standard fiscal shocks. Their simulations demonstrate that oil-financed investment booms produce strong short-run growth, whereas output contracts sharply when oil revenue declines and investment is cut. They also emphasize that a stabilization fund could substantially reduce output volatility. These results illustrate the broader risk: fiscal consolidation strategies are difficult to evaluate in commodity exporters because outcomes hinge on factors outside the direct control of policymakers.

Cross-country evidence reinforces this view. \citet{VillafuerteLopez-MurphyOssowski13} document how fiscal policies in Latin American resource exporters follow a ``roller-coaster'' dynamic, expanding during booms and retrenching during busts, amplifying the cycle and complicating consolidation efforts. \citet{vanPloegVenables2011} and \citet{Venables16} stress that windfall revenues are often mismanaged, with governments failing to build buffers that would support adjustment in downturns. Forecasting studies confirm that commodity prices are hard to predict, with large errors even at short horizons \citep[see][]{BaumeisterKilian2016}, which means that baseline revenue projections used in consolidation programs are subject to high risk.

One response has been the development of fiscal frameworks that delink spending from current commodity revenue. Chile’s structural balance rule is the most prominent example, where expert panels estimate long-run copper prices to set expenditure paths \citep[see][]{Frankel2011}. Empirical evidence shows that resource exporters with fiscal rules are less procyclical and less vulnerable to revenue-driven retrenchments \citep[see][]{Bova2014,Eyraud2018}. While these studies focus on improving fiscal stability more broadly, they do not directly assess how commodity revenue volatility interacts with fiscal consolidation plans. Our paper addresses this gap by explicitly modeling the vulnerability of consolidation strategies to commodity price shocks, using a structural framework that captures the dynamic interplay between fiscal instruments, external conditions, and macroeconomic outcomes.

A complementary empirical literature studies how commodity terms-of-trade shocks propagate through small open economies and interact with macro policy. For example, \citet{Roch17} provides a comparative analysis of Chile, Colombia, and Peru, documenting how exchange-rate regimes and fiscal space shape the adjustment to commodity price declines. Related evidence for advanced commodity exporters highlights sectoral reallocation effects: \citet{KnopVespignani14} show that commodity price shocks have heterogeneous impacts across Australian industries, with mining and construction responding strongly relative to other sectors.\footnote{For related evidence on heterogeneous transmission in Australia, see \citet{Vespignani15}, who documents differential monetary-policy effects across states/territories in a small open economy setting.} For oil exporters, \citet{AlJabrietal22} document that oil price shocks account for a substantial share of fiscal revenue and output variation in Oman, and emphasize fiscal smoothing through buffers and debt issuance. Relative to this empirical small-open-economy evidence on commodity shocks, our contribution is to provide a structural DSGE laboratory that maps a real-world, multi-instrument consolidation program into explicit policy paths and quantifies how plausible commodity-revenue deviations translate into debt and spread dynamics during implementation.

From a modeling perspective, DSGE frameworks for commodity exporters have formalized key macro-fiscal dynamics. \citet{MedinaSoto2016} and \citet{MelinaYangZanna2016} show how commodity price shocks interact with fiscal rules and debt dynamics, emphasizing the importance of institutional design in stabilizing resource-rich economies. \citet{Buffie2012} embed public investment and debt sustainability in a DSGE model tailored to low-income, resource-rich countries, illustrating how the financing and composition of fiscal adjustment shape growth outcomes. While these studies focus on the broader macroeconomic implications of commodity shocks and fiscal policy, they do not explicitly examine how revenue volatility affects the credibility or effectiveness of fiscal consolidation plans.

Finally, the broader fiscal consolidation literature warns that composition and timing influence the real effects of adjustment. Spending-based consolidations are often less damaging in advanced economies \citep[see][]{AlesinaFaveroGiavazzi2015}, but in commodity exporters, where capital spending is tightly linked to commodity receipts, cutting investment can be particularly contractionary \citep[see][]{Guajardo2014}. In Latin America, evidence suggests that consolidation through reduced public investment has especially large multipliers, making such strategies both painful and risky when driven by revenue shortfalls \citep[see][]{CarriereSwallow2018}.

Taken together, this literature highlights the challenges of designing and evaluating fiscal consolidations in commodity-exporting economies. Our contribution is to formalize and quantify the vulnerability of consolidation plans to plausible deviations in commodity revenue forecasts, using a structural small open economy model with explicit fiscal and commodity blocks. While we ground our simulations in the case of Ecuador, the framework and insights are broadly applicable to other resource-dependent economies facing similar fiscal risks.

%%%%%%%%%%%%%%%%%%%%%%%%%%%%%%%%%%%%%%%%%%%%%%%%%%%%%%

\section{The DSGE Model}

%%%%%%%%%%%%%%%%%%%%%%%%%%%%%%%%%%%%%%%%%%%%%%%%%%%%%%

This section spells out a two-agent SOE model with real rigidities and a rich fiscal policy block. The economy is populated by Ricardian households that have access to international financial assets, and rule-of-thumb households that consume their disposable income every period in hand-to-mouth fashion. There are three categories of government expenditure: (i) consumption of final goods, which enters household utility as a proxy for publicly provided goods and services (e.g., health and education), (ii) investment in public capital, which enters the production function of the firms producing domestic goods, and (iii) transfers to households. The government finances its expenditures by levying three types of distortionary taxes\textemdash on consumption, labor income, and capital income\textemdash as well as by issuing foreign debt. Domestic goods, which are produced with private and public capital in addition to labor, are combined with imported goods to form final goods. Real rigidities include habit formation in private consumption and adjustment costs in private investment. The economy trades in domestic and foreign goods and exports an exogenous endowment of a commodity. Finally, the economy is subject to shocks to preferences, technology (neutral and investment-specific), commodity production, the components of government expenditure, tax rates, foreign GDP, foreign interest rates, a risk premium on external debt, and the international price of the commodity good.

{\normalsize \label{sec:model}}

\subsection{Households}

There is a continuum of infinitely lived households of unit mass. A fraction $\omega\in (0,1)$ of households are restricted and do not have access to saving instruments. They behave in a hand-to-mouth fashion whereby they consume their disposable income every period; we use the superscript $hm$ to denote variables that correspond to them. The remaining fraction $(1-\omega)$ of households have access to financial assets and can accumulate physical capital, which they rent to firms; we refer to this group as Ricardian households and denote variables related to them with the superscript $R$.\footnote{This specification follows the standard two-agent structure widely used in DSGE models with fiscal policy \citep[e.g.][]{coenen2005does, gali2007understanding}.} Both types of households have preferences defined over a composite consumption bundle, $\hat{C}^j_{t}$, that includes both private and public components, and hours worked, $h^j_{t}$, for $j\in\{R,hm\}$, that are paid a common hourly wage, $W_t$.\footnote{Throughout, uppercase letters denote variables containing a unit root in equilibrium (due to technological progress), whereas lowercase letters indicate variables with no unit root. The imported good is the numeraire. Appendix \ref{Appendix:equations} describes how each variable is transformed to achieve stationarity in equilibrium. Variables without a time subscript denote non-stochastic steady state values.}

Following \citet{coenen2013gauging}, the consumption bundle for households of type $j=\{R,hm\}$ is a constant-elasticity-of-substitution (CES) aggregate of household purchases, denoted $\hat{C}_t^j$, of a final good, $C_t^j$, and government consumption, $G_t^c$, as follows:

\begin{equation*}
\hat{C}_t^j \equiv \hat{C}^j(C_t^j,\breve{C}_{t-1}^j,G_t^c)= \left[(1-o_c)^{\frac{1}{\eta_c}}\left(C_t^j-\varsigma \breve{C}_{t-1}^j\right)^{\frac{\eta_c-1}{\eta_c}}+o_c^{\frac{1}{\eta_c}}\left(G_t^c\right)^{\frac{\eta_c-1}{\eta_c}}\right]^{\frac{\eta_c}{\eta_c-1}},
\end{equation*}
where $\breve{C}_t^j$ denotes average consumption purchases across households of type $j$, which each household takes as given (in equilibrium, $ C_t^j=\breve{C}_t^j$). The parameter $\varsigma\in (0,1)$ governs consumption habits, whereas $\eta_c>0$ governs the elasticity of substitution between consumption purchases by the household and the government.\footnote{$\eta_c \to 0$ implies perfect complementarity, $\eta_c \to \infty$ gives perfect substitutability, and $\eta_c \to 1$ yields the Cobb–Douglas case.}

Expected discounted utility of a representative household of type $j$ is given by
\begin{equation}
E_{t}\sum_{s=0}^{\infty }\beta ^{s}v_{t+s}\left[ \frac{1}{1-\sigma }\left(\hat{C}^j_{t+s}\right) ^{1-\sigma }-\Theta^j_{t+s}A_{t+s-1}^{1-\sigma }\kappa _{t+s}^j\frac{\left(h^j_{t+s}\right)^{1+\phi }}{1+\phi }\right], 
\label{eq:Households_1}
\end{equation}%
where $\upsilon_t$ and $\kappa _{t}$ are exogenous preference shocks, $\beta \in (0,1)$ is the intertemporal discount factor, $\sigma >0$ and $\phi \geq 0$ are the inverse elasticity of intertemporal substitution and the inverse Frisch elasticity of hours worked, respectively, whereas $A_{t}$ is a non-stationary technology index. Following \citet{gali2012unemployment}, $\Theta^j_{t}$ is an endogenous preference shifter that regulates the strength of the wealth effect on labor supply.\footnote{The preference shifter $\Theta^j_{t}$ satisfies:
\begin{equation*}
\Theta^j_{t}= \widetilde{\chi }_{t}^j A_{t-1}^{\sigma }\left[\hat{C}^j(C_t^j,\breve{C}_{t-1}^j,G^c)\right]^{-\sigma },\qquad \widetilde{\chi }_{t}^j=\widetilde{\chi }%
_{t-1}^{j^{1-\nu }}A_{t-1}^{-\sigma \nu }\left[\hat{C}^j(\breve{C}_t^j,\breve{C}_{t-1}^j,G^c)\right] ^{\sigma \nu },
\end{equation*}
where $\nu \in \left[ 0,1 \right]$ regulates the strength of the wealth effect: when $\nu=1$, preferences are of the standard CRRA type; when $\nu=0$, the wealth effect is shut down, as in the preferences by \citet{greenwood1988investment}. Because the consumption bundle that determines household utility includes government consumption, the preference shifter $\Theta^j_t$ includes steady-state government consumption $G^c$, but this is a proxy. We acknowledge that the actual variable has a unit root and is transformed accordingly for stationarity.}

\paragraph{Ricardian Households.}

Ricardian households save and borrow by trading non-state-contingent international bonds, $F_{t}^{R\ast }$, with foreign agents (foreign variables will be denoted with an asterisk). They also purchase an investment good, $I^R_{t}$, which determines next period's physical capital stock, $K^R_{t}$. Further, let $r_{t}^{\ast }$ denote the gross return on $F_{t-1}^{R\ast }$ and let $r_t^k$ be the rental rate on capital $K_{t-1}^R$. Additionally, let $\Pi_t^R$ denote the profits from the home good producing firm, which is owned by this type of household.

Ricardian households pay taxes on labor income, consumption expenditures, and capital income (after deducting depreciation costs) at rates, $\tau^n_t, \tau^c_t$, and $\tau^k_t$, respectively. They also receive government transfers, $TR^R_t$. Hence, the period-by-period budget constraint of the representative Ricardian household is given by%
\begin{equation}
(1+\tau^c_t)p_tC^R_{t}+p_tI^R_{t} +r_{t}^{\ast}F^{R\ast}_{t-1} + T^R 
=(1-\tau^n_t)W_{t}h^R_{t}+F^{R*}_t+ \left[(1-\tau^k_t)r_{t}^{K}+\tau^k_t\delta \right]K^R_{t-1} + TR^R_t + \Pi_t^R, %+r_{t-1}B^R_{t-1}
\label{eq:Households_3}
\end{equation}%
where $p_t$ denotes the price of the final good relative to the price of the imported good, which is the numeraire.\footnote{Note that consumption, $C_t^R$, and investment, $I_t^R$, are bundles of the final good, so to express them in terms of the foreign good, they must be multiplied by $p_t$; the other terms in the budget constraint are expressed in terms of the foreign good.} $T^R$ is a constant lump-sum tax introduced to ensure consistency between the model’s steady-state budget constraint and observed fiscal aggregates, such as the share of government consumption and transfers in GDP.

The physical capital stock evolves according to the following law of motion:

\begin{equation}
K^R_{t}=(1-\delta )K^R_{t-1}+[1-\Gamma \left( I^R_{t}/I^R_{t-1}\right)]u_{t}I^R_{t},\qquad \delta \in [0,1],  \label{eq:Households_4}
\end{equation}%
where $I^R_{t}$\ denotes investment expenditures and%
\begin{equation*}
\Gamma \left( \frac{I^R_{t}}{I^R_{t-1}}\right) =\frac{\gamma_k}{2}\left( \frac{I^R_{t}}{I^R_{t-1}}-\bar{a}\right) ^{2},\qquad \gamma \geq 0,\qquad \bar{a}\geq
1,
\end{equation*}%
are convex investment adjustment costs. The variable $u_{t}$ is an investment shock that captures changes in the efficiency of the investment process \citep[see][]{justiniano2010investment} and $\bar{a}$ is the steady-state trend growth rate of the economy. The household chooses $C^R_{t}$, $h^R_{t}$, $I^R_{t}$, $K^R_{t}$, and $F^{R\ast}_{t}$ to maximize (\ref{eq:Households_1}) subject to (\ref{eq:Households_3})-(\ref{eq:Households_4}), taking $W_{t}$, $r_{t}^{\ast }$, $r_{t}^{K}$, $TR^R_{t}$, $G^c_{t}$, $\Pi_t^R$, $F^{R\ast}_{t-1}$, $K^R_{t-1}$, and $T^R$ as given.  

The (gross) interest rate at which Ricardian households save and borrow is given by:%
\begin{equation}
r_{t}^{\ast } =R_{t-1}^{\ast }\xi _{t-1}, \label{eq:interestrate}
\end{equation}
where $R_{t-1}^{\ast }$ is the gross foreign real interest rate, which evolves exogenously, and $\xi _{t-1}$ denotes a country premium given by \citep[see][]{schmitt2003closing, Adolfson2008evaluating}:
\begin{equation}
\label{eq:country_premium}
\xi _{t}=\bar{\xi}\exp \left[ \psi \frac{(F^{\ast}_t +B_{t}^{\ast})/Y_{t}-(\bar{f}^{\ast} + \bar{b}^{\ast})}{\bar{f}^{\ast} + \bar{b}^{\ast}}+\frac{\zeta _{t}-\bar{\zeta} }{%
\bar{\zeta} }\right] ,\qquad \psi >0,\qquad \bar{\xi}\geq 1,
\end{equation}%
where $F^{\ast}_t$ is the aggregate stock of private foreign debt, $B_t^*$ is foreign government debt, $Y_t$ is GDP (to be defined below), $\zeta _{t}$ is an exogenous idiosyncratic component of the country premium, and $(\bar{f}^{\ast}+\bar{b}^{\ast}$) is an exogenous target for the ratio of foreign debt to GDP.\footnote{Both the government and Ricardian households borrow at the same interest rate, given by the world interest rate plus a risk premium that depends on the economy’s total external-debt-to-GDP ratio. This specification implies that changes in external indebtedness affect the borrowing costs of both the public and private sectors. Ecuador is officially dollarized, which effectively fixes the nominal exchange rate whereas capital mobility is imperfect in the model, so domestic borrowing costs adjust through the sovereign risk premium or the world interest rate rather than through domestic monetary policy or exchange-rate movements.} The steady-state level of the country risk premium is determined by $\bar{\xi}$. The country premium increases with the economy’s external debt burden and with exogenous risk shocks, capturing the idea that a higher indebtedness level and adverse global conditions raise the economy's borrowing costs.\footnote{The model represents government liabilities as a single debt instrument priced at the external interest rate plus a sovereign risk premium. In practice, a country can issue both domestic and external public debt. A standard alternative formulation in SOE DSGE models assumes that a fixed fraction $\alpha \in (0,1)$ of public liabilities is held abroad while the remaining share is held domestically. Under integrated capital markets and common pricing conditions, this decomposition is largely notational and does not affect the equilibrium dynamics of fiscal consolidation or the link between public debt and sovereign spreads. We therefore interpret the government debt variable in the model, $B_t^*$, as one of the marginal borrowing instruments relevant for determining the country's risk premium. As a reference, external debt represented around 70\% of Ecuador’s public debt during the sample period.}

Our model determines the real interest rate and country premium without an explicit sovereign default decision. While Ecuador’s 2008 debt restructuring lies within our estimation sample (2004–2019), we treat it as a sharp repricing of sovereign risk\textemdash captured by large estimated spread shocks\textemdash rather than as an endogenous default event. This approach is standard in SOE DSGE models where the focus is on the macroeconomic transmission of borrowing costs rather than the strategic default decision.\footnote{Introducing a full sovereign-default block would require solving an occasionally binding, dynamic problem in which the government endogenously chooses whether to repay or default, creditors update expectations about future restructuring terms, and equilibrium bond prices reflect these dynamics. Implementing such a framework is feasible but significantly enlarges the state space, introduces non-differentiabilities, and generates multiplicity issues. Our main contribution\textemdash evaluating how alternative commodity-revenue paths shape the macroeconomic and fiscal consequences of a consolidation plan\textemdash operates through debt dynamics and sovereign spreads rather than through the extensive-margin default decision. A tractable intermediate option would be to embed a reduced‑form sovereign‑risk block, in which default risk rises mechanically with debt or projected deficits, as in \citet{Corsettietal13} or \citet{Cantoreetal19}. These frameworks link sovereign spreads to fiscal fundamentals without modeling strategic default, offering a middle ground between our current Schmitt‑Grohé–Uribe structure and full default models. Although appealing in some contexts, incorporating such a mechanism does not materially affect the channels we study and would therefore not alter our main findings.}

\paragraph{Hand-to-Mouth Households.}

The subset of households that do not have access to financial assets or capital accumulation receive only labor income and government transfers, $TR^{hm}_t$, which they use to consume and pay the consumption tax. We assume that hand-to-mouth households do not pay labor income taxes, reflecting either thresholds in the tax code or their representation of informal sector workers. Hence, they face the following simpler budget constraint:

\begin{equation}
\label{eq:Households_5}
(1+\tau^c_t)p_tC^{hm}_{t}=W_{t}h^{hm}_{t}+ TR^{hm}_t.
\end{equation}%
They solve a static optimization problem, choosing $C_t^{hm}$ and $h_t^{hm}$ to maximize \eqref{eq:Households_1}, subject to \eqref{eq:Households_5}.

\subsection{Firms}

The supply side of the economy is composed of four types of competitive firms, as follows: a set of firms that produce a home good using labor and capital, a set of firms that import a homogeneous foreign good, a set of firms that combine the domestic and foreign goods into a final good that is purchased by households and the government for consumption and investment, and a set of competitive firms producing a homogeneous commodity good that is exported. A proportion of those commodity-exporting firms is owned by the government and the remaining proportion is owned by foreign agents. The total mass of firms in each sector is normalized to one.

\paragraph{Home Goods.}

Home goods are produced according to the following technology:%
\begin{equation}
Y_{t}^{H}=z_{t}\left(K^g_{t-1}\right)^{\gamma}K_{t-1}^{\alpha }[A_{t}h_{t}]^{1-\alpha - \gamma },\qquad
\alpha \in \lbrack 0,1), \qquad \gamma \in \lbrack 0,1), \label{eq:IntermGoods_1}
\end{equation}%
where $z_{t}$ is an exogenous stationary technology shock, $A_{t}$ (with $a_{t}\equiv A_{t}/A_{t-1}$) is a non-stationary technology index, $K_{t-1}$ is the aggregate stock of private capital available for production in period $t$, and $K^g_{t-1}$ is the stock of public capital.
The firm hires labor and rents private capital to maximize profits subject to the technology constraint \eqref{eq:IntermGoods_1}, taking the stock of public capital as given.\footnote{The firm earns positive profits in equilibrium due to the inclusion of public capital, which introduces decreasing returns to scale in private inputs, thereby breaking the standard zero-profit condition under constant returns.}

\paragraph{Foreign Goods.}

A representative importing firm buys an amount $Y_t^F$ of a homogeneous foreign good to meet the demand of the final good producers, $X_t^F$. This representative importing firm makes zero profits.

\paragraph{Final Goods.}

A\ representative final-goods firm demands home and foreign goods in the amounts $X_{t}^{H}$ and $X_{t}^{F}$, respectively, and combines them according to the following technology:%
\begin{equation}
Y_{t}^{C}=\left[ \left( 1-o\right) ^{\frac{1}{\eta }}\left( X_{t}^{H}\right)
^{\frac{\eta -1}{\eta }}+o^{\frac{1}{\eta }}\left( X_{t}^{F}\right) ^{\frac{%
\eta -1}{\eta }}\right] ^{\frac{\eta }{\eta -1}},\qquad o\in (0,1),\qquad
\eta >0.  \label{eq:FinalGoods_1}
\end{equation}%

Let $p_t$ and $p_{t}^{H}$ denote the relative prices of $Y_t^C$ and $X_{t}^{H}$, respectively, in terms of the imported good. Subject to the technology constraint (\ref{eq:FinalGoods_1}), the firm maximizes its profits $\Pi_{t}^{C}=p_tY_{t}^{C}-p_{t}^{H}X_{t}^{H}-X_{t}^{F}$ over the input demands $X_{t}^{H}$ and $X_{t}^{F}$, taking $p_t$ and $p_{t}^{H}$ as given. This firm makes zero profits in equilibrium.

\paragraph{Commodities.}

A representative commodity-producing firm produces a quantity of a commodity good $Y_{t}^{P}$\ in each period, which evolves exogenously. The entire production is sold abroad at a given international price, $p_{t}^{P*}$, which also evolves exogenously. The  income generated in the commodity sector is therefore equal to $p_{t}^{P*}Y_{t}^{P}$. The government receives a share $\chi \in \lbrack 0,1]$ of this income and the
remaining share goes to foreign agents.

We model commodity output, $Y_t^P$, as an exogenous stochastic process. This choice is motivated by the low-frequency and highly persistent dynamics of Ecuadorian oil production over our estimation window, and by the fact that short-run fluctuations in oil output are weakly co-moving with non-oil activity. Our focus is therefore on the macro-fiscal transmission of commodity \emph{revenue} shocks---driven primarily by international prices and exogenous output movements---through the fiscal block, external debt, and the country risk premium, rather than on the determinants of oil-sector investment, depletion, and operational constraints.\footnote{Endogenizing commodity production would require introducing an explicit oil sector with capital accumulation, investment and adjustment costs, and possibly depletion or infrastructure constraints. Such an extension is feasible but would move the focus away from our main objective: quantifying how uncertainty in commodity revenues alters the performance of a given fiscal consolidation plan.}

\subsection{Fiscal Policy}

The government demands the final good for consumption ($G^c_{t}$) and investment ($G^i_{t}$) purposes and makes transfers to households ($TR_t$). It levies taxes on consumption expenditures, labor and capital income, issues foreign debt, and receives a share $\chi\in(0,1)$ of the income generated in the commodity sector. Hence, the government satisfies the following period-by-period budget constraint:\footnote{The tax rates $\tau^c_t$, $\tau^n_t$ and $\tau^k_t$ are fixed at their steady state values when we estimate the model, but in the simulation of a fiscal consolidation, which includes increases in tax rates, we allow them to follow the paths specified in the IMF consolidation plan.}

\begin{equation}
p_tG^c_{t} + p_tG^i_{t} + TR_{t}  + r_t^{\ast}B_{t-1}^{\ast} = B_t^* + \tau^c_tp_tC_t + (1-\omega)\tau^n_tW_th_t^R + \tau^k_t(r^k_t-\delta)K_{t-1} + \chi p_{t}^{P*}Y_{t}^{P} + (1-\omega)T^R. %+  r_{t-1}B_{t-1} %+ B_t
\label{eq:FiscalPolicy}
\end{equation}%

Government expenditure components respond endogenously to the output gap (deviations of output from steady state) and deviations of the government debt-to-GDP ratio from target, capturing both cyclical behavior and potential fiscal sustainability concerns, as follows:

\begin{equation}
G_t^c = (a_{t-1}G_{t-1}^c)^{\rho_{gc}}\left(G^c \left(\frac{Y_t}{Y} \right)^{\alpha_{gc}} \left(\frac{B^{\ast}_{t-1}/Y_{t-1}}{\bar{b}^*} \right)^{\gamma_{gc}}\right)^{1-\rho_{gc}}\exp(e_t^{gc}), \label{eq:gc}
\end{equation}
\begin{equation}
G_t^i = (a_{t-1}G_{t-1}^i)^{\rho_{gi}}\left(G^i \left(\frac{Y_t}{Y} \right)^{\alpha_{gi}} \left(\frac{B^{\ast}_{t-1}/Y_{t-1}}{\bar{b}^*} \right)^{\gamma_{gi}}\right)^{1-\rho_{gi}}\exp(e_t^{gi}), \label{eq:gi}
\end{equation}
\begin{equation}
TR_t = (a_{t-1}TR_{t-1})^{\rho_{TR}}\left(TR \left(\frac{Y_t}{Y} \right)^{\alpha_{TR}} \left(\frac{B^{\ast}_{t-1}/Y_{t-1}}{\bar{b}^*} \right)^{\gamma_{TR}}\right)^{1-\rho_{TR}}\exp(e_t^{TR}), \label{eq:tr}
\end{equation}
where $\rho_{gc}$, $\rho_{gi}$ and $\rho_{TR}$ lie in the $[0, 1)$ interval and govern the persistence of the components of expenditure. The parameters $\alpha_{gc}$, $\alpha_{gi}$, and $\alpha_{TR}$ govern the strength of the response of each expenditure component to the business cycle, with positive values indicating procyclical behavior and negative values indicating countercyclical behavior. The parameters $\gamma_{gc}$, $\gamma_{gi}$, and $\gamma_{TR}$ govern the strength of the response of each expenditure component to deviations of the government debt-to-GDP ratio from its target, $\bar{b}^*$, with negative values indicating a debt-stabilizing behavior. The components of government expenditure are also subject to exogenous white noise shocks, $e_t^{gc}$, $e_t^{gi}$, and $e_t^{TR}$.\footnote{Once again, we are slightly abusing notation in the expenditure rules by denoting the steady state value of each component as $G^c$, $G^i$ and $TR$. The components of expenditure have a unit root, so they do not have a steady state. It is the transformed variables, $\frac{G_t^c}{A_{t-1}}$, $\frac{G_t^i}{A_{t-1}}$, and $\frac{TR_t}{A_{t-1}}$ that have a steady state.}

These expenditure reaction functions are intended to summarize observed fiscal behavior in the pre-consolidation sample through two standard margins: cyclical comovement (via the output gap term) and fiscal sustainability (via feedback from the debt-to-GDP gap). This type of specification is widely used in DSGE models for emerging and commodity-exporting economies as a tractable way to ensure public debt remains bounded while allowing for empirically relevant procyclicality in spending components \citep[see, e.g.,][]{MedinaSoto2016,MelinaYangZanna2016,Buffie2012}. The estimated coefficients therefore inform the model’s baseline stochastic dynamics and help discipline the joint Bayesian estimation of structural and policy parameters. Compared to structural-balance frameworks that delink spending from contemporaneous commodity revenues by using long-run commodity price benchmarks \citep[e.g., Chile’s rule discussed in][]{Frankel2011}, our specification is a reduced-form representation of observed behavior in Ecuador over the estimation sample. This choice reflects our focus on evaluating a given consolidation plan under commodity-revenue uncertainty rather than on designing an optimal fiscal anchor.

Government investment augments the stock of public capital, which is an input in the production of home goods, according to the following law of motion:
\begin{equation*}
K^g_{t} = (1-\delta_g)K^g_{t-1} + G^i_{t}, \quad \delta_g \in [0,1],
\end{equation*}
where $\delta_g$ is the depreciation rate of public capital.

The government distributes shares $\omega^G$ and $(1-\omega^G)$ of total transfers to hand-to-mouth and Ricardian households, respectively. Total transfers received by hand-to-mouth households satisfy $\omega TR_t^{hm}=\omega^G TR_t$, whereas transfers received by Ricardian households satisfy $(1-\omega)TR_t^{R}=(1-\omega^G) TR_t$. This formulation allows the government to target transfers differently than the population share, for example focusing them on the financially restricted hand-to-mouth households.

The IMF’s Extended Fund Facility provides official loans to the government at concessional rates. For tractability, we do not distinguish official from market borrowing in the model. External financing conditions are summarized by the country‑risk‑premium block, and our consolidation simulations condition on the fiscal‑instrument paths specified by the program\textemdash not on the composition or pricing of the government’s financing sources. Introducing a separate official‑lending class would primarily alter the composition of external liabilities and marginal interest expenditures, but it would not affect the dynamic macroeconomic responses we highlight.

\subsection{The Rest of the World}

Foreign demand for the home good, $X_{t}^{H\ast }$, is given by the schedule%
\begin{equation*}
X_{t}^{H\ast }=o^{\ast }\left(p_t^H\right)^{-\eta^*}Y_{t}^{\ast },\qquad o^{\ast }\in (0,1), \qquad \eta\ast >0,
\end{equation*}%
where $Y_{t}^{\ast }$ denotes foreign aggregate demand, proxied by foreign GDP, and $\eta^*$ is an elasticity that governs how exports demand responds to changes in the price of the home good.\footnote{Detrended foreign output is modeled as an exogenous AR(1) process that captures fluctuations in external demand. It is independent of the domestic productivity process that drives the trend growth of the economy. This assumption reflects the small open economy perspective adopted in the model, whereby domestic developments do not influence foreign output.}

\subsection{Aggregation and Market Clearing}

To aggregate private consumption, hours worked, and transfers, consider that both types of households, Ricardian and hand-to-mouth, consume and work, and receive transfers. Hence, the following conditions hold: $C_t=(1-\omega)C_t^R+\omega C_t^{hm}$; $h_t=(1-\omega)h_t^R+\omega h_t^{hm}$; and $TR_t=(1-\omega)TR_t^R+\omega TR_t^{hm}$.

As only Ricardian households own firms and hold financial assets and physical capital, we have the following conditions for aggregate profits, private asset holdings, the aggregate capital stock and aggregate investment, respectively: $\Pi_t= (1-\omega)\Pi_t^R$; $F^{\ast}_t=(1-\omega)F_t^{R\ast}$; $K_t=(1-\omega)K_t^R$; and $I_t=(1-\omega)I_t^R$.

In the market for final goods, the clearing condition is $Y_{t}^{C}=C_{t}+I_{t}+G^c_{t} +G^i_{t}$, while in the markets for home and foreign goods, we have $Y_{t}^{H}=X_{t}^{H}+X_{t}^{H\ast }$, and $Y_{t}^{F}=X_{t}^{F}$. Hence, aggregate demand is the sum of domestic absorption, $Y_{t}^{C}$, and the trade balance, $TB_t$, which in equilibrium equals real GDP, $Y_t$. Consequently, expressed in terms of foreign goods (the numeraire), real GDP is given by:
\begin{equation*}
Y_{t} = p_tY_t^C+TB_t=\underbrace{p_t(C_{t}+I_{t}+G^c_{t} +G^i_{t})}_{\text{Domestic absorption}}+\underbrace{p_t^H X_{t}^{H\ast }+p_t^{P\ast}Y_{t}^{P}-Y_{t}^F}_{\text{Trade balance}}.
\end{equation*}
Alternatively, real GDP can be expressed as domestic supply: $Y_t = p_t^H Y_t^H + p_{t}^{P\ast }Y_{t}^{P}$.

The net foreign asset position of this economy results from combining the period-by-period budget constraints of the households and government. The aggregate household budget constraint is the weighted sum of the Ricardian and hand-to-mouth households' budget constraints, where the weights are their shares in the population, $(1-\omega)$ and $\omega$, respectively, which results in:
\begin{align*}
(1+\tau^c_t)p_tC_{t}+p_tI_{t}+r^*_tF^{\ast}_{t-1} + (1-\omega)T^R =&W_{t}h_{t} -(1-\omega)\tau^n_tW_{t}h_{t}^R+F^{\ast}_{t} +  \\
& +\left[(1-\tau^k_t)r_{t}^{K}+\tau^k_t\delta \right]K_{t-1} + TR_t + \Pi_t. %+r_{t-1}B_{t-1} %+B_{t}
\end{align*}
Combining the previous expression with the period-by-period budget constraint of the government in \eqref{eq:FiscalPolicy} obtains the balance-of-payments identity:
\begin{equation*}
F^{\ast}_t + B_{t}^{\ast }=r_{t}^{\ast }(F^{\ast}_{t-1}+ B_{t-1}^{\ast})-TB_{t}+(1-\chi)p_{t}^{P\ast }Y_{t}^{P}, %
\end{equation*}
which states that the change in the net foreign asset position is equal to the current account balance (the trade balance plus net interest income).

\subsection{Driving Forces}

In addition to the previously described shocks to tax rates and government expenditures, several exogenous variables in the model follow an AR(1) process of the form
\begin{equation*}
    log(x_t/\bar{x})=\rho_x log (x_{t-1}/\bar{x})+\varepsilon_t^x, \qquad \rho_x\in [0,1),
\end{equation*}
\sloppy
where $\varepsilon_t^x$ are white noise shocks and $\bar{x}$ is their steady-state counterpart, for $x=\{v, \kappa, u, z, a, \zeta, R^{\ast}, p^{P*}, y^{P}, y^{\ast}\}$. The economy is subject to a wide array of shocks, including those to consumption and labor supply preferences, investment efficiency, technology (both transitory and permanent), external financial conditions, and commodity market dynamics. These shocks reflect the key sources of macroeconomic volatility faced by commodity-exporting countries.\footnote{The full set of stationary equilibrium conditions, including the first-order conditions associated with the optimization problems described above appear in Appendix \ref{Appendix:equations}.}

In sum, the model integrates heterogeneous households, a detailed fiscal block, and a realistic external sector to capture the macroeconomic dynamics of a commodity-exporting small open economy. Its structure allows us to analyze how fiscal consolidation plans interact with commodity revenue volatility, providing a rich framework for policy evaluation.

%%%%%%%%%%%%%%%%%%%%%%%%%%%%%%%%%%%%%%%%%%%%%%%%%%%%%%

\section{Parametrization Strategy}\label{sec:parametrization}

Our strategy for taking the model to the data combines the estimation of most of the parameters and the calibration of the rest. While some parameters are directly assigned values based on external sources, particularly studies focused on Latin American economies, others are solved endogenously within the steady-state system to match empirical targets, such as ratios with respect to GDP of government consumption, transfers, external debt, and oil production, and the public-to-total investment ratio.\footnote{See Appendix \ref{Appendix:steadystate} for details on the parameters that are solved endogenously.}

\paragraph{Calibrated Parameters.}
In general, our calibrations for preferences and technology follow the strategies in \citet{MedinaSoto2016}, whose model is tailored to Latin American economies with similar structural features, and \citet{coenen2013gauging}. To obtain targets for some of the steady-state values, we use the 2022-2027 projections in the July 2022 IMF Staff Report for Ecuador \citep[see][]{oner2022imf} and data counterparts over the estimation sample in other cases.\footnote{Table \ref{table:steadystatecalibration} in Appendix \ref{app:calibratedparameters} presents the calibrated parameters and targeted steady state values, and their sources.}

\paragraph{Estimation of Parameters.}
We estimate the parameters of the model with Bayesian techniques, solving the model with a linear approximation around the non-stochastic steady state. Priors are selected based on standard values in the literature and economic plausibility, ensuring convergence of the posterior distributions. We use quarterly information from Ecuador for the period 2004:Q1\textendash 2019:Q4 for the following 13 observable variables: Non-oil GDP, private aggregate consumption, aggregate investment, government consumption, government investment, government transfers, the trade balance, hours worked, the country risk premium as measured by the Emerging Markets Bond Index (EMBI) for Ecuador, the West Texas Intermediate (WTI) oil price, oil production, the foreign interest rate, and foreign GDP. The inclusion of oil prices and the EMBI spread helps identify parameters related to commodity shocks and the country risk premium. We exclude post-2019 data mainly to ignore pandemic-related distortions.\footnote{Table \ref{table:priorposterior} in Appendix \ref{app:estimatedparameters} shows the prior and posterior distributions of the estimated parameters. In the interest of keeping the discussion in the main body of the paper fluid and concise, this appendix also discusses the parameter estimates.}

We now turn to evaluating the model’s ability to replicate key empirical moments and its dynamic behavior under the estimated parameters.


\section{Model Dynamics and Performance}\label{sec:performance}

To assess how the model fits the data, we compute for nine selected variables the following model-implied moments, evaluated at the posterior mean of the parameters: their correlation with GDP, standard deviation, and first-order autocorrelation coefficient. We then compare these statistics with empirical moments. Overall, the model replicates key empirical features of the Ecuadorian economy, including the co-movement with GDP, volatility, and persistence of macroeconomic aggregates, with deviations generally within acceptable bounds for estimated DSGE models.\footnote{Table \ref{table:data_model_moments} in Appendix \ref{app:moments} reports the comparison between model-implied and empirical moments.}

To understand how the model interprets the data through its structural shocks, Figure \ref{figure:shock_decomp_gdp} shows the historical shock decomposition of the annual growth rate of non-oil GDP.\footnote{Table \ref{table:vardecomp} in Appendix \ref{app:moments} reports the unconditional variance decomposition of selected variables.} Shocks are grouped into six categories: i) commodity shocks, which include oil price ($p^{P*}$) and oil output ($Y^{P}$) shocks, ii) a risk shock that affects the country premium ($\zeta$), iii) fiscal policy shocks, which include those to government consumption, government investment, and transfers ($e^{gc}$, $e^{gi}$, and $e^{TR}$), iv) technology shocks, which include transitory ($z$) and permanent ($a$) productivity shocks, and investment efficiency ($u$) shocks, v) foreign shocks, which include shocks to the foreign interest rate ($R^*$) and foreign demand ($y^*$), and vi) preference shocks, which include consumption ($v$) and labor supply ($\kappa$) shocks.\footnote{Historical decompositions in DSGE models should be interpreted as statistical allocations of fluctuations across structural disturbances rather than literal identifications of specific historical events. We ignore the role of initial conditions and measurement error, so the bars do not add up exactly to the solid line. Appendix \ref{Appendix:shockdecomp} contains the shock decomposition of the other observable variables of the model.} 

Ecuador displays a cycle marked by a substantial expansion since the mid 2000s, which, as in other Latin American countries, coincided with the boom in international commodity prices. This boom was briefly interrupted in 2008\textendash 2009 by the Global Financial Crisis (GFC), and later brought to a sudden halt in 2014, when international commodity prices collapsed. The last five years of the sample are marked by a seemingly permanent decline in the growth rate of non-oil GDP.\footnote{Looking at the non-oil GDP growth rate through the lens of the model points to technology shocks (dark blue bars) as the most important drivers. This is particularly relevant in the last part of the sample, in which permanent productivity shocks ($a_t$) are the only mechanism through which the model can rationalize the persistent decline in output growth. Because the productivity process contains a permanent component that determines the long-run growth rate of non-oil output, the estimation attributes structural changes in trend growth to this disturbance. Other shocks in the model---including fiscal, oil, risk-premium, and preference shocks---are stationary and therefore cannot generate permanent shifts in the growth rate of output. As in any estimated DSGE model, these disturbances should be interpreted as reduced-form residuals capturing persistent forces not explicitly modeled (e.g., institutional changes, sectoral reallocation, or external demand shifts), rather than as literal productivity innovations. A deeper structural explanation of this slowdown is beyond the scope of this paper. One possible interpretation is that growth patterns featuring weak (strong) demand induce a slowdown (rise) in productivity, as argued by \citet{stadler1990business}, \citet{anzoategui2019endogenous}, \citet{jorda2020long}, and \citet{furlanetto2025estimating}, among others.}

One of the objectives of this paper is to understand the roles that oil prices, country risk, and fiscal policy play in the fiscal consolidation proposed in the IMF program with Ecuador. As an initial approach to understand these roles, we analyze the historical shock decomposition of output growth fluctuations in our sample. Oil-related shocks (gray bars) are relevant drivers of output growth during episodes of oil price swings, such as the high prices observed from the mid 2000s to 2013, which exert upward pressure, in general. Oil shocks contribute negatively during the GFC, when the oil price declined, and especially during 2014\textendash 2017, a period of low oil prices that exerted substantial negative pressure on output growth. 

\begin{figure}[ht]
    \centering
     \caption{Historical Shock Decomposition of Annual Non-Oil GDP Growth}
    \includegraphics [width=\textwidth]{Figures/gama_YNCo_obs.eps}
    \label{figure:shock_decomp_gdp}
    \smallskip
    \captionsetup{width=.8\linewidth}
    \caption*{\footnotesize{\emph{Note}: The solid black line denotes the four-quarter moving sum of the observable variable\textemdash the demeaned quarter-on-quarter log difference of non-oil GDP, which approximates its annual growth rate. The bars are four-quarter moving sums of shocks obtained from the historical shock decomposition of the observable variable. For details on the shock groupings, see the text.}}
\end{figure}

Country risk premium shocks (yellow bars) exert downward pressure on output growth around the GFC of 2008 for two reasons, one global and one idiosyncratic. First, the global turmoil increased international investors' risk perception of emerging markets. The second and surely more important reason is that Ecuador decided to default on its sovereign debt for political considerations in December of 2008. These forces sent the Ecuadorian EMBI spread above 4,000 points (a 40 percentage-point premium above the equivalent risk-free bond). As time progressed and Ecuador regained access to financial markets, the downward trend in the country risk premium generated a positive effect of these shocks on output growth through 2014. Between 2015 and 2016, risk shocks put downward pressure on growth, again due to global and idiosyncratic factors. In this period, country premiums increased in developing countries (especially in Latin America) due to the decline in commodity prices. The idiosyncratic factor was a severe earthquake that affected two coastal provinces in Ecuador. Towards the end of the sample, with the country risk premium mostly declining, the effect of these shocks on GDP growth is small but positive.

Finally, fiscal policy shocks (light blue bars) contribute positively to output growth during 2007\textendash 2008, a period of highly expansionary fiscal policy. Fiscal policy shocks are also expansionary during 2016\textendash 2017, likely due to the severe earthquake and its effect on the Ecuadorian economy, which required a fiscal expansion in full force to restore physical infrastructure.\footnote{Governmental and international institutions report the cost of reconstruction and economic recovery was close to \$3 billion. See \url{https://www.worldbank.org/en/news/feature/2021/04/27/a-cinco-a-os-del-terremoto-ecuador-sigue-trabajando-en-su-resiliencia-frente-a-desastres} and \url{https://ecuadorchequea.com/moreno-ley-de-solidaridad-1500-millones/}.} Additionally, the presidential elections of 2017 might have also played a role in the expansionary fiscal policy, as the sitting president was trying to promote a candidate of his own party.\footnote{It is worth noting that the relatively small contribution of foreign shocks reflects their estimated historical variance rather than their theoretical importance in the model. During the sample period, Ecuador’s business cycle was dominated by large oil-price fluctuations, fiscal adjustments, and movements in the sovereign risk premium. As a result, the estimation attributes most of the variation in domestic output to these shocks rather than to innovations in foreign output. Similar findings are common in estimated DSGE models of commodity-exporting economies, where terms-of-trade and domestic fiscal shocks typically account for a large share of macroeconomic fluctuations \citep[see][for instance]{MedinaSoto2016, MelinaYangZanna2016}.}

To illustrate how two shocks of interest for this paper\textemdash oil price and country risk premium shocks\textemdash propagate through the model economy, we analyze their impulse-response functions. These shocks are interesting because they may affect fiscal outcomes and macroeconomic dynamics during the implementation of fiscal consolidation plans, potentially generating deviations from baseline trajectories.

%%%%%%%%%%%%%%%%%%%%%%%%%%%%%%%%%%%%%%%%%%%%%%%%%%%%%%

\paragraph{Effects of an oil price shock.}
A one-standard-deviation shock to the price of oil causes the price to increase by 19\% on impact, as shown in figure \ref{irf_oil} (top-left panel). The government receives an oil revenue windfall that %The higher oil price leads to more revenue for the government, which drives down 
reduces the deficit (panel 2-1) and, accordingly, public debt (panel 2-2). The decline in the model economy's debt-to-GDP ratio leads to a decline of more than 20 bps in the country premium (panel 2-4). Following the estimated reaction functions of government expenditure, all three components\textemdash government consumption, government investment, and transfers\textemdash increase gradually, reaching peaks four to twelve quarters after the shock hits the economy (panels 1-2, 1-3, and 1-4). The expansion in government expenditure reflects both improved fiscal space\textemdash via a lower debt-to-GDP ratio\textemdash and the endogenous estimated procyclical response of spending to a positive output gap, as encoded in the fiscal rules.

\begin{figure}[ht]
    \centering
     \caption{Impulse Responses to an Oil Price Shock}
    \includegraphics [scale=.75]{Figures/oil_irf.eps}
    \label{irf_oil}
    \smallskip
    \captionsetup{width=.6\linewidth}
    \caption*{\footnotesize{\emph{Note}: Impulse responses to a one-standard-deviation shock to the international oil price. X axes denote quarters. Y axes denote percent deviations from steady state, except for variables expressed as ratios to GDP and the country premium, which are expressed as percentage point deviations from steady state. Public and private debt, and the country premium, are annualized. }}
\end{figure}

Because public and private consumption are complementary for household utility, both the Ricardian and hand-to-mouth households increase consumption on impact (panels 4-1, 4-2, and 3-2). The increase in transfers also contributes to the increase in consumption of hand-to-mouth households, who also reduce hours worked (panel 4-4) due to the income effect. The decline in the country premium increases the present value of investment---i.e., increases Tobin's q---so the Ricardian households increase private investment (panel 3-3), such that the 20 bps decline in the country premium translates into a roughly 0.4\% increase in investment a year after the shock hits the economy. The overall effect on non-oil GDP is slightly negative on impact (panel 3-1) because hours worked reach their trough almost immediately whereas public and private investment take several quarters to reach their peak. Non-oil GDP begins expanding persistently about a year after the shock hits the economy.

In summary, an increase in oil prices raises government revenue, improving the fiscal balance and reducing the need for external borrowing. Lower external debt leads to a decline in the sovereign risk premium, which reduces the domestic interest rate faced by both the government and private agents. This stimulates private investment and consumption. Through these channels, oil-price movements propagate to the external and fiscal balances and, ultimately, to output.

\paragraph{Effects of a country premium shock.}
We now describe the response of the model economy to a shock to the country risk premium, which is shown in figure \ref{irf_risk}. A one-standard-deviation shock to the idiosyncratic component of the country premium (the top-left panel) drives up the premium by over 300 basis points (panel 2-4). This increase in the country premium drives down the present value of investment, which leads the Ricardian households to decrease investment (panel 3-3). The higher country premium also exerts downward pressure on the consumption of Ricardian households (panel 4-1) due to  the intertemporal substitution effect. Furthermore, because Ricardian households are debtors in steady state, the higher country premium makes them relatively poorer, which also drives down their consumption, as well as their leisure (panel 4-3), due to the negative income effect. 

The higher country premium also increases the government's interest payments and, thus, the deficit (panel 2-1) and debt (panel 2-2). The components of government expenditure (panels 1-2, 1-3 and 1-4) decline for two reasons\textemdash to stabilize a rising public debt-to-GDP ratio and due to their procyclical response to a contractionary shock. The decline in government transfers depresses consumption of the hand-to-mouth households (panel 4-2) and leads them to work more (panel 4-4) due to the negative income effect. Overall, non-oil GDP rises slightly on impact due to the expansion of hours worked, but declines persistently shortly after the shock hits the economy due to the fall in investment.

Eventually, the decline in government expenditure pushes the deficit-to-GDP ratio below steady state (panel 2-1), containing the increase in public debt. Forward-looking, Ricardian households anticipate that this fiscal contraction will induce the country premium to fall below steady state, which ultimately pushes their consumption and investment above steady state. 

\begin{figure}[ht]
    \centering
    \caption{Impulse Responses to an Idiosyncratic Country Risk Shock}
    \includegraphics [scale=.75]{Figures/risk_irf.eps}
    \label{irf_risk}
    \smallskip
    \captionsetup{width=.6\linewidth}
    \caption*{\footnotesize{\emph{Note}: Impulse responses to a one-standard-deviation shock to the idiosyncratic component of country risk. Y axes denote percent deviations from steady state, except for variables expressed as ratios to GDP and the country premium, which are expressed as percentage point deviations from steady state. Public and private debt, and the country premium, are annualized.}}
\end{figure}

Building on these dynamics, we now examine how fiscal consolidation plans perform under alternative paths for commodity revenue, highlighting the vulnerabilities introduced by external shocks.

\section{Fiscal Consolidations} \label{sec:fiscalpolicy}

This section discusses the macroeconomic effects of fiscal consolidations under historically plausible deviations in commodity revenue from forecasts. To ground our simulations, we use the fiscal consolidation agreement between the government of Ecuador and the IMF, which we used in section \ref{sec:parametrization} to discipline our model (henceforth the agreement), for the 2020\textendash 2025 period. We analyze the model economy's performance when oil revenue is higher or lower than originally envisioned in the program.\footnote{Our simulation approximates the consolidation program, recognizing that the actual fiscal adjustments involve institutional and distributional complexities beyond the scope of our stylized DSGE framework.}

The agreement between Ecuador and the IMF consisted of a combination of increases in tax rates and a reduction in government expenditure. We conduct the simulations under perfect foresight: agents observe the full quarter‑by‑quarter trajectories of tax rates, government expenditure, and oil revenue over the 24‑quarter consolidation horizon. During this period, these variables follow the paths specified in the IMF‑supported program and we impose them directly as deterministic sequences, rather than being generated by the fiscal‑policy reaction functions estimated for the pre‑program sample.\footnote{The estimated reaction functions for government expenditure (\ref{eq:gc})\textendash (\ref{eq:tr}) are used to match historical co-movements and to identify policy parameters jointly with the rest of the model. The consolidation experiment, by construction, conditions on the program’s expenditure trajectories; thus the reaction functions mainly affect the estimated dynamics, not the consolidation counterfactual.} After the consolidation window, we assume that tax rates remain quasi‑permanently at their new levels, the components of government expenditure revert to their estimated reaction functions (\ref{eq:gc})\textendash (\ref{eq:tr}), and oil revenue gradually returns to its steady state at the rate implied by the estimated persistence of the oil‑price process.\footnote{After the consolidation window, we assume that tax rates evolve as follows: $\tau^x_t = .99\tau^x_{t-1} + (1-.99)\tilde{\tau}^x + \varepsilon _{\tau^x}$, for $x\in \{c,n,k\}$ where $\tilde{\tau}^x$ is the steady-state tax rate.}

The trajectories of tax rates and government expenditures used to generate our simulation are shown in tables \ref{tIMF_taxes} and \ref{tIMF_spending}, respectively. Tax rate changes, which are expressed as percentage point (pp) changes from tax rates in the previous period, include a permanent 3 pp increase in the labor income tax rate in 2022, the third year of the program, and a transitory 0.5 pp increase in the capital income tax rate during the third year of the program. The agreement did not include changes to the value-added tax rate, but they could be implemented as deterministic trajectories of the consumption tax rate, as described above.\footnote{The personal income tax reform implied a more progressive scheme and a reduction of deductions for high-income earners. Because the model has a single labor income tax rate for all agents, we assume an average tax rate change across the board as a way to implement this reform, taking into account that the Ecuadorian government estimated that 3.4\% of the labor force would have been affected by the reform (see the proposal of the Law for Economic Development and Fiscal Sustainability sent by the Executive branch to Congress: \url{https://www.comunicacion.gob.ec/wp-content/uploads/2021/11/LeyParaElDesarrolloEconomico.pdf}). Similarly, the consolidation program considered a temporary contribution in 2022 by persons whose net worth was above \$1 million (a 1\% to 1.5\% tax rate) and firms whose net worth was above \$5 million (0.8\% tax rate). To include these measures in our model, we assume a transitory increase in the capital income tax rate. For more details, see IMF Country Report No. 22/225 (\url{https://www.elibrary.imf.org/view/journals/002/2022/225/article-A000-en.xml}).}

Regarding government expenditure, table \ref{tIMF_spending} shows the evolution of its three components as percent changes from the previous year in real terms. Broadly speaking, the consolidation program contemplates a moderation of the growth rate of government consumption, a substantial decline in the level of government investment in the first year that is only partially recovered in the following years, and a trajectory for the level of transfers that consists of a decline in the first year, a substantial increase in the second and third years, and another decline for the final three years. The increase in transfers was designed to protect vulnerable households from adverse effects of the consolidation program.\footnote{The consolidation simulations condition on the fiscal‑instrument paths specified by the IMF‑supported program and do not impose any path for official disbursements. External financing is captured in reduced form through the country‑risk‑premium block, and public debt evolves endogenously given the fiscal stance. As a result, the simulations abstract from the composition of external government borrowing (official versus market) and focus on the macro‑fiscal implications of the program’s policy instruments.}

\begin{table}[ht] 
      \centering 
     \caption{Fiscal Consolidation: Tax Rate Changes}  \label{tIMF_taxes} %IMF Consolidation:
     \begin{threeparttable}
     \begin{tabular}{l d{2.1} d{2.1} d{2.1} d{2.1}  }
  \hline \hline
\multicolumn{1}{c}{Year}& \multicolumn{1}{c}{Consumption Tax}   & \multicolumn{1}{c}{Labor Income Tax}     & \multicolumn{1}{c}{Capital Income Tax}  \\  
\hline \hline
2020& +0 & +0 & +0	 \\
2021&+0 &+0 &+0 \\
2022&+0 &+3 &+0.5\\
2023& +0 &+0 &-0.5\\
2024&+0 &+0 &+0\\
2025&+0 &+0 &+0\\
\hline
 \end{tabular}
 \begin{tablenotes}
\small
\item \footnotesize \emph{Note}:  Entries denote percentage point changes from tax rates in the previous year.
\end{tablenotes}
     \end{threeparttable}
\end{table}

\begin{table}[ht] 
      \centering 
     \caption{Fiscal Consolidation: Government Expenditure Trajectory}  \label{tIMF_spending}
     \begin{threeparttable}
     \begin{tabular}{l d{2.1} d{2.1} d{2.1}  }
  \hline \hline
\multicolumn{1}{c}{Year}& \multicolumn{1}{c}{Government Consumption}  & \multicolumn{1}{c}{Government Investment}   & \multicolumn{1}{c}{Transfers}  \\  
\hline \hline
2020& +0.9 & -29.4 & -25.7  \\
2021& +1.8 & +10.6 & +57.9  \\
2022& -3.9 & -1.6 & +11.1  \\
2023& +0.8 & +4.1 & -14.2  \\
2024& +0.5 & +0.9 & -14.8  \\
2025& +1.9 & +0.8 & -15.0  \\
\hline
 \end{tabular}
 \begin{tablenotes}
\small
\item  \footnotesize \emph{Note}: Entries denote percent changes from the previous year in real terms.
\end{tablenotes}
     \end{threeparttable}
\end{table}

A final ingredient in the simulation is the trajectory of oil revenue. Table \ref{tab:tIMF_oilrevenue} shows the projected real growth rate of oil revenue in the consolidation agreement. Since the program was approved on September 2020, about six months after the onset of the COVID-19 pandemic, it projected a substantial decline in the level of oil revenue in that year, a strong recovery in the following two years, and a moderation in the last three years of the program.

\begin{table}[ht] 
      \centering 
     \caption{Fiscal Consolidation: Oil Revenue Trajectory}  \label{tab:tIMF_oilrevenue}
     \begin{threeparttable}
     \begin{tabular}{c c c c c c }
  \hline \hline
2020& 2021& 2022& 2023& 2024& 2025 \\
\hline \hline
-42.3 & +63.9 & +30.4 & -9.5 & -9.4 & -8.6 \\
\hline
 \end{tabular}
 \begin{tablenotes}
\small
\item  \footnotesize \emph{Note}: Entries denote percent changes from the previous year in real terms.
\end{tablenotes}
     \end{threeparttable}
\end{table}

A high-level description of the effects of fiscal consolidation is as follows (a detailed discussion appears below). The effects arise from the interaction between household decisions, production, and financial conditions. Fiscal instruments affect households through their budget constraints and intratemporal incentives: higher labor taxes distort labor supply by reducing after-tax wages, while transfers directly affect the income and labor-supply choices of hand-to-mouth households. Public consumption influences household utility, particularly when it is complementary to private consumption, thereby affecting consumption–labor trade-offs. Public investment operates through the accumulation of public capital, which enters firms’ production functions and affects the marginal productivity of private inputs. At the same time, fiscal consolidation reduces the need for external borrowing, putting downward pressure on the sovereign risk premium. The resulting decline in borrowing costs lowers the user cost of capital and stimulates private investment. These effects are heterogeneous across households: hand-to-mouth households respond primarily to contemporaneous income changes, while Ricardian households adjust consumption, labor supply, and investment intertemporally. The overall dynamics of the economy reflect the interaction between labor-supply adjustments, the evolution of private and public capital, and the financial channel operating through the risk premium along the transition path.

Figure \ref{fig:imf} shows the macroeconomic effects of this consolidation plan under the baseline oil revenue path and under scenarios with positive and negative deviations from the forecast. Solid blue lines display the baseline case, in which the real growth rate of oil revenue follows the projected path in the program. Dashed orange lines show an alternative case in which oil revenue is one historical standard deviation higher, whereas oil revenue is one standard deviation lower in the scenario in dotted yellow lines.\footnote{The trajectories of the six \textit{fiscal instruments} is the same in all scenarios. This is obvious for the case of tax rates, shown in panels 2-1 and 2-2 of figure \ref{fig:imf}. The level of each of the components of government expenditure follows the same trajectory in real terms across the three scenarios, but when expressed as ratios to GDP, as in panels 1-1, 1-2 and 1-3, they differ. This is due to differences in the denominator (GDP), which is an endogenous outcome of each simulation and is not immune to alternative paths for oil revenue. The standard deviation of oil revenue is computed from the historical quarterly growth rate of real oil revenue over the estimation sample. We then add/subtract this value to/from the baseline growth rate in each quarter, preserving the original trajectory’s shape but shifting its level.}

\begin{figure}[ht]
    \centering
    \caption{Fiscal Consolidation Under Alternative Paths for Oil Revenue}
    \label{fig:imf}
    \includegraphics [scale=.8]{Figures/IMF_Simulation.eps}
    \smallskip
    \captionsetup{width=.8\linewidth}
    \caption*{\footnotesize{\emph{Note}: Perfect foresight simulations of a fiscal consolidation combining adjustments of multiple revenue and expenditure instruments designed to approximate the agreement between the government of Ecuador and the IMF, as shown in tables \ref{tIMF_taxes} and \ref{tIMF_spending}. The simulations also consider paths for oil revenue. Solid blue lines depict the effects of the fiscal consolidation with oil revenue as forecast by the agreement (see table \ref{tab:tIMF_oilrevenue}). The dashed orange and dotted yellow lines display cases when oil revenue is one standard deviation higher and lower, respectively, than expected by the agreement. X axes denote quarters. Y axes denote percent deviations from steady state, except for variables expressed as ratios to GDP and the country premium, which are expressed as percentage point deviations from steady state. The country premium is annualized.}}
\end{figure}

We first comment on the baseline case, describing the trajectories of the fiscal instruments and their macroeconomic effects (the blue lines). Consistent with the values shown in table \ref{tIMF_taxes}, the labor income tax rate permanently increases by 3 pp from the third year (quarter 9, panel 2-1), the capital income tax rate increases by 0.5 pp during the third year (panel 2-2), and the consumption tax rate (not shown) remains at steady state over the entire simulation. The annual paths of the components of government spending experience more variability, as shown in table \ref{tIMF_spending}. To get sequences of quarterly pp deviations of GDP ratios from steady state, we apply a quadratic interpolation to the annual sequences, ensuring that each year's four-quarter average is equal to the corresponding annual value. After the initial increase we described before, government consumption declines gradually, reaching a trough of slightly more than 1 pp of GDP below steady state in quarter 12 (top-left panel). Government investment declines immediately, with the government investment-to-GDP ratio fluctuating around 0.7 pp below steady state throughout the 24-quarter horizon (panel 1-2). The transfers-to-GDP ratio displays an inverse U-shape, lying below steady state during the first year, rising above it during the second to fourth years, and dipping below for the final two years (panel 1-3). Finally, in line with the values in table \ref{tab:tIMF_oilrevenue}, the path of the oil revenue-to-GDP ratio in panel 1-4 starts about 2 pp below steady state during the first year, then rises gradually for the next two years, reaching a peak of about 1 pp in quarter 12 of the simulation, and then declines in the remaining three years, reaching about -1 pp by the end of the horizon.

The fiscal consolidation reduces the deficit-to-GDP ratio (panel 2-3), which reaches 2 pp below steady state by quarter 12, and 3 pp below steady state by the end of the horizon.\footnote{In our model, the effect of the consolidation on the fiscal balance is smaller than the IMF estimation; Country Report No. 22/225 states that the primary balance of the non financial public sector would improve by 4.5\% of GDP between 2020 and 2025 with the program.} The decline in the deficit (panel 2-3) contains public debt, but private debt (not shown) increases to finance private consumption and investment, which initially puts upward pressure on the premium. The eventual decline in overall debt leads to a decline of the country risk premium, reaching about 35 bps below steady state by the end of the consolidation horizon (panel 2-4), which in turn increases private investment (panel 3-3) because of the lower cost of capital.

Higher income taxes reduce the consumption of Ricardian households (panel 4-1), whereas the consumption of restricted, hand-to-mouth households (panel 4-2) follows the inverse U-shape pattern of government transfers, as this group of consumers receives the entirety of transfers. The inverse U-shaped response of aggregate consumption (panel 3-2) reflects the interaction of short-run and medium-run effects. In the short run, higher labor taxes and lower transfers reduce disposable income, leading to an overall decline in consumption. At the same time, lower public expenditure reduces overall absorption. Over time, the fiscal consolidation reduces external debt and lowers the sovereign risk premium, which decreases borrowing costs and stimulates private investment and the consumption of Ricardian households. As capital accumulates and income rises, consumption gradually recovers, especially for liquidity-constrained households. Ricardian households also smooth consumption intertemporally, contributing to the recovery.

Labor-supply responses differ across households due to their access to financial markets. Ricardian households adjust hours worked based on both intratemporal substitution (changes in after-tax wages) and intertemporal considerations, as they can smooth consumption over time. Hand-to-mouth households, in contrast, do not participate in financial markets and therefore choose labor supply based only on current income conditions, responding to intratemporal incentives. Both types of households are also affected by wealth effects: higher income or transfers tend to reduce labor supply, while lower income increases it. As a result, fiscal consolidation—through higher labor taxes and lower transfers—tends to increase labor supply, particularly among hand-to-mouth households (panel 4-4), while Ricardian households exhibit a less pronounced response due to their intertemporal optimization (panel 4-3).

Taken together, the dynamics of hours worked (lower) and capital, both public (lower) and private (higher), imply a decline in non-oil GDP of approximately 0.5\% below trend during years 2 and 3 (panel 3-1). From the fourth year, the expansion in the capital stock through its private component contribute to a recovery of GDP, which ends the six-year simulation horizon slightly below trend.\footnote{The trajectory of the trade balance (panel 3-4) is largely driven by the path of oil revenues. Changes in domestic absorption—reflecting the interaction between private and public consumption and investment—also affect imports, but their net contribution is more muted relative to the direct effect of oil revenues.}

We now turn to the alternative scenarios, where oil revenue deviates positively or negatively from the baseline projection, to assess how such shocks affect the consolidation’s effectiveness.\footnote{We treat commodity production as exogenous; hence cross‑scenario differences arise from commodity revenues (prices and exogenous production) rather than endogenous supply responses. Introducing an explicit oil sector with investment/capacity constraints could amplify persistence in revenue deviations, but our central result would remain: under fixed fiscal‑instrument paths, revenue realizations map into different debt/spread trajectories and macroeconomic outcomes.} The alternative scenarios for oil revenue generate substantially different macroeconomic outcomes, holding the trajectories in tax rates and government expenditure fixed. The path of the oil revenue-to-GDP ratio (panel 1-4) is roughly 2 pp higher or lower under historically plausible deviations. When oil revenue is higher, as depicted by dashed orange lines, the deficit-to-GDP ratio (panel 2-3) declines about twice as much as in the baseline case, reaching 5 pp below steady state by the end of the horizon. The lower deficit contains public debt such that the decline in the country risk premium reaches more than 100 bp (panel 2-4). In these conditions, private investment expands forcefully (panel 3-3). Real non-oil GDP (panel 3-1) barely dips below trend and begins increasing in the second year.

When oil revenue is lower than anticipated, as depicted by dotted yellow lines, the deficit-to-GDP ratio (panel 1-4) barely declines below steady state. Without further adjustment to tax rates or government expenditure, this situation generates an increase in the country risk premium, which reaches about 50 bp above steady state by the end of the horizon (panel 2-4). Private investment (panel 3-3) declines in these circumstances, while private consumption (panel 3-2) declines more than in the baseline simulation. Real non-oil GDP (panel 3-1) experiences a larger and more persistent decline, remaining more than 1\% below trend by the end of the simulation horizon.\footnote{As an additional sensitivity exercise, we explore how the baseline fiscal consolidation would change when the pass-through from a lower risk premium is not complete, as would be the case, for example, in a country that imposes capital controls. The results appear in Appendix \ref{app:capitalcontrols}.}

The perfect-foresight simulations assume the announced program is credible. We acknowledge that Ecuador’s historical fiscal slippages and episodes of policy reversal imply that perfect credibility is an unrealistic assumption. Limited credibility would amplify the risks associated with low oil revenue paths: if agents expected possible future deviations from the announced fiscal plan, the risk premium would incorporate an additional credibility spread, worsening the macro-fiscal dynamics.\footnote{A full formal treatment of imperfect credibility would require a Markov-switching or occasionally-binding policy regime, which is beyond the scope of this paper, but would be an important extension in future work.}

To summarize this section: The consolidation package combines several fiscal instruments that affect the economy through real and financial channels. Increases in labor-income taxes distort labor supply decisions by reducing after-tax wages and disposable income. Changes in capital taxation affect investment decisions through their impact on the user cost of capital. Reductions in transfers directly affect hand-to-mouth households, altering their income and labor-supply choices. Cuts in public consumption affect household utility and intratemporal allocations, without directly entering production. Reductions in public investment have more persistent effects, as they slow the accumulation of public capital and thus reduce future productive capacity. In addition to these direct effects, fiscal consolidation reduces the need for external public borrowing, putting downward pressure on the sovereign risk premium. The resulting decline in borrowing costs lowers the user cost of capital and stimulates private investment, supporting the accumulation of private capital. The overall dynamics of the consolidation reflect the interaction between labor-supply adjustments, capital accumulation—both private and public—and the financial channel operating through the risk premium. Appendix \ref{appendix:impulseresponses} illustrates the effects of adjusting each fiscal instrument separately to achieve a fiscal deficit reduction of one percent of GDP.\footnote{To quantify parameter uncertainty around these transition paths, we report posterior‑draw credible intervals in Appendix~\ref{app:uncertainty}.}

\section{Conclusion}\label{sec:conclusion}

This paper explored the macroeconomic effects of fiscal consolidations in commodity-exporting countries using a DSGE model estimated for Ecuador and used to simulate its 2020\textendash 2025 IMF-supported adjustment program. Our simulations show that the outcomes of fiscal consolidation plans can be highly sensitive to deviations in commodity revenue from forecasts.

When oil revenue is higher than expected, the fiscal deficit improves more rapidly, public debt declines, and the country risk premium falls. These dynamics may reduce the perceived urgency of continued fiscal adjustment and could, in practice, create incentives to relax or abandon the consolidation effort. Conversely, when oil revenue falls short of projections, the fiscal deficit deteriorates compared with the original projection, the country risk premium rises, and private investment weakens. In such cases, achieving the original fiscal targets may require a more aggressive adjustment, which could be politically or socially difficult to implement.

More broadly, the paper contributes to the literature on fiscal policy in commodity-exporting small open economies in three ways. First, we evaluate a multi-instrument IMF-style fiscal consolidation using a fully specified structural SOE DSGE model, allowing us to trace nonlinear interactions among commodity revenue realizations, external debt dynamics, and sovereign risk premia under perfect foresight. Second, the model highlights the macro-financial mechanism through which commodity-revenue fluctuations interact with external debt dynamics and sovereign spreads, shaping the evolution of output, investment, and the fiscal balance during consolidation. Third, by incorporating heterogeneous households, the framework clarifies how fiscal adjustments affect different groups through distinct channels, with hand-to-mouth households responding mainly through labor income and transfers and Ricardian households adjusting consumption and investment intertemporally.

While our model does not explicitly capture the political economy of fiscal policy, the results suggest that the design of consolidation programs should account for the risks posed by commodity price volatility. Incorporating mechanisms such as flexible fiscal targets or contingency clauses could help programs remain credible and sustainable in the face of external shocks, ensuring that temporary revenue windfalls or shortfalls do not derail the broader adjustment effort.

Finally, our quantitative results are disciplined by Ecuador’s data and institutional context, but the core mechanism applies more broadly to commodity exporters: for a given consolidation plan, commodity‑revenue realizations shift debt and sovereign spreads, which propagate to macro outcomes through financing conditions. The strength of this channel will vary with institutions. In economies with monetary autonomy and flexible exchange rates, exchange‑rate adjustment and monetary policy can partially buffer adverse commodity shocks, attenuating (though not eliminating) spread amplification. Countries with stronger non‑commodity tax bases, stabilization funds, or structural‑balance frameworks can smooth fiscal adjustment and reduce sensitivity to near‑term commodity revenues. Conversely, economies closer to fiscal limits or facing steeper debt–spread schedules may exhibit stronger nonlinear amplification, making consolidation outcomes more fragile to revenue shortfalls. We therefore view the magnitudes reported here as Ecuador‑specific, but the qualitative lesson—consolidation performance is highly contingent on commodity‑revenue realizations and the financing conditions they induce—is relevant across a wide set of resource‑dependent economies.

\clearpage
\newpage
\bibliographystyle{jmcb}
\singlespacing
\bibliography{references.bib}

\clearpage
\appendix
\appendixpage
\numberwithin{equation}{section}
\numberwithin{table}{section}
\numberwithin{figure}{section}

\section{Stationary Equilibrium Conditions}
\label{Appendix:equations}


\renewcommand{\theequation}{A.\arabic{equation}} \setcounter{equation}{0}

The variables in uppercase contain a unit root in equilibrium due to the presence of the non-stationary productivity index $A_{t}$. We need to transform these variables to have a stationary version of the model. To do this, lowercase variables denote the uppercase variable divided by $A_{t-1}$ (e.g. $%
c_{t}\equiv \frac{C_{t}}{A_{t-1}}$). The only exception is the Lagrange multiplier $\Lambda _{t}$, which is multiplied by $A_{t-1}^{\sigma }$ (i.e. $\lambda _{t}\equiv \Lambda _{t}A_{t-1}^{\sigma }$), for it decreases along
the balanced growth path.

The rational expectations equilibrium of the stationary version of the model
model is the set of 45 sequences
\begin{equation*}
\{\lambda_{t}^R,\hat{c}_t^R,c_{t}^R,h_{t}^R,i_{t}^R,k_t^R, f_t^{*o},\Theta_t^R,\widetilde{\chi}_t^R, \lambda_{t}^{hm},
\hat{c}_t^{hm},c_{t}^{hm},h_{t}^{hm}, \Theta_t^{hm}, \widetilde{\chi}_t^{hm},c_{t},h_{t},i_{t},k_t,f^*_t,q_t, w_t, r_t^{K},y_{t},y_{t}^{C},y_{t}^{F}\}
\end{equation*}
\begin{equation*}
\{y_{t}^{H},x_{t}^{F},x_{t}^{H},x_{t}^{H\ast},\pi_t,p_{t}^{H},p_{t},tb_{t},\xi_t,b^*_t,\tau^k_t,\tau^n_t,\tau^c_t,k^g_t,g^c_t,g^i_t,tr_t,tr^R_t,tr^{hm}_t\},
\end{equation*} such that for given initial values and 11 exogenous sequences%
\begin{equation*}
\{v_t,\kappa^R _{t}, \kappa^{hm}_t,u_{t},z_{t},a_{t},\zeta _{t},R_{t}^{\ast },p_{t}^{Co\ast },y_{t}^{Co},y_{t}^{\ast }\}
\end{equation*}%
the following conditions are satisfied:%

\noindent Ricardian consumption first order condition:
\begin{equation}
\lambda _{t}^R(1+\tau_t^c)p_t=\hat{c}_t^{o^{-\sigma}}\left(\frac{(1-o_c)\hat{c}^R_t}{c_t^R - \varsigma\frac{c_{t-1}^R}{a_{t-1}} }\right)^{\frac{1}{\eta_c}} \label{eq:EE_1}
\end{equation}%
Hand-to-mouth consumption first order condition:
\begin{equation}
\lambda _{t}^{hm}(1+\tau_t^c)p_t=\hat{c}_t^{r^{-\sigma}}\left(\frac{(1-o_c)\hat{c}^{hm}_t}{c_t^{hm} - \varsigma\frac{c_{t-1}^{hm}}{a_{t-1}} }\right)^{\frac{1}{\eta_c}} \label{eq:EE_2}
\end{equation}%
Consumption aggregation:
\begin{equation}
c_t = (1-\omega)c_t^R + \omega c_t^{hm}
\label{eq:EE_3}
\end{equation}
Ricardian hours first order condition:
\begin{equation}
(1-\tau_t^n)w_{t}=\Theta _{t}^R\kappa^R _{t}\frac{h_{t}^{o^\phi }}{%
\lambda^R _{t}}  \label{eq:EE_4}
\end{equation}%
Hand-to-mouth hours first order condition:
\begin{equation}
w_{t}=\Theta _{t}^{hm}\kappa^{hm} _{t}\frac{h_{t}^{r^\phi }}{%
\lambda^{hm} _{t}}  \label{eq:EE_5}
\end{equation}%
Hours aggregation:
\begin{equation}
h_t = (1-\omega)h_t^R + \omega h_t^{hm} \label{eq:EE_6}
\end{equation}
Ricardian debt first order condition:
\begin{equation}
\lambda^R _{t}=\frac{\beta }{a_{t}^{\sigma }}R_{t}^{\ast }\xi _{t}E_{t}\left\{
\frac{v_{t+1}}{v_{t}}\lambda^R _{t+1}\right\}
,  \label{eq:EE_7}
\end{equation}%
Final good production function:
\begin{equation}
y_{t}^{C}=\left[ \left( 1-o\right) ^{\frac{1}{\eta }}\left( x_{t}^{H}\right)
^{\frac{\eta -1}{\eta }}+o^{^{\frac{1}{\eta }}}\left( x_{t}^{F}\right) ^{%
\frac{\eta -1}{\eta }}\right] ^{\frac{\eta }{\eta -1}},  \label{eq:EE_8}
\end{equation}%
Foreign good first order condition:
\begin{equation}
x_{t}^{F}=o\left(\frac{1}{p_t}\right) ^{-\eta }y_{t}^{C},  \label{eq:EE_9}
\end{equation}%
Home good first order condition:
\begin{equation}
x_{t}^{H}=\left( 1-o\right) \left(\frac{p_{t}^{H}}{p_{t}}\right) ^{-\eta }y_{t}^{C},
\label{eq:EE_10}
\end{equation}%
Foreign demand for home goods:
\begin{equation}
x_{t}^{H\ast }=o^{\ast }\left( p_{t}^{H}\right) ^{-\eta
^{\ast }}y_{t}^{\ast },  \label{eq:EE_11}
\end{equation}%
Home good production function:
\begin{equation}
y_{t}^{H}=z_{t} \left(\frac{k_{gt-1}}{a_{t-1}}\right)^{\gamma} \left(\frac{k_{t-1}}{a_{t-1}}\right) ^{\alpha
}(a_{t}h_{t})^{1-\alpha-\gamma },  \label{eq:EE_12}
\end{equation}%
Home good market clearing:
\begin{equation}
y_{t}^{H}=x_{t}^{H}+x_{t}^{H\ast },  \label{eq:EE_13}
\end{equation}%
Final good market clearing:
\begin{equation}
y_{t}^{C}=c_{t}+i_{t}+g^c_{t}+g^i_{t},  \label{eq:EE_14}
\end{equation}%
Defenition of GDP:
\begin{equation}
y_{t}=p_t\left(c_{t}+i_{t}+g^c_{t}+g^i_{t}\right)+tb_{t}=p_t^H y_t^H+p_t^{Co\ast}y_{t}^{Co},  \label{eq:EE_15}
\end{equation}%
Defenition of trade balance:
\begin{equation}
tb_{t}=p_{t}^{H}x_{t}^{H\ast }+p_{t}^{Co\ast }y_{t}^{Co}-y_{t}^F,
\label{eq:EE_16}
\end{equation}%
Private debt aggregation:
\begin{equation}
f_t^{\ast} = (1-\omega)f_t^{o\ast} \label{eq:EE_17}
\end{equation}
Balance of payments:
\begin{equation}
(f_t^* + b_{t}^{\ast })=\frac{(f_{t-1}^*+b_{t-1}^{\ast })}{a_{t-1}}R_{t-1}^{\ast }\xi _{t-1}-tb_{t}+(1-\chi )p_{t}^{Co\ast }y_{t}^{Co},
\label{eq:EE_18}
\end{equation}%
Country premium:
\begin{equation}
\xi _{t}=\bar{\xi}\exp \left[ \psi \frac{(f_t^* + b_{t}^{\ast})/y-(\bar{f}^{\ast} + \bar{b}^{\ast})}{\bar{f}^{\ast} + \bar{b}^{\ast}}+\frac{\zeta _{t}-\bar{\zeta} }{\bar{\zeta} }\right] ,
\label{eq:EE_19}
\end{equation}%
Evolution of private capital:
\begin{equation}
k_{t}^R=(1-\delta )\frac{k_{t-1}^R}{a_{t-1}}+\left[ 1-\frac{\gamma_k }{2}\left(
\frac{i^R_{t}}{i^R_{t-1}}a_{t-1}-\bar{a}\right) ^{2}\right] u_{t}i^R_{t},
\label{eq:EE_20}
\end{equation}%
Private capital aggregation:
\begin{equation}
k_t = (1-\omega)k_t^R
\label{eq:EE_21}
\end{equation}
Private investment aggregation:
\begin{equation}
i_t = (1-\omega)i_t^R
\label{eq:EE_22}
\end{equation}
Ricardian capital first order condition:
\begin{equation}
q_{t}=\frac{\beta }{a_{t}^{\sigma }}E_{t}\left\{ \frac{v_{t+1}}{v_{t}}\frac{%
\lambda _{t+1}^R}{\lambda _{t}^R}\left[ (1-\tau^k_{t+1})r_{t+1}^{K}+\tau^k_{t+1}\delta + q_{t+1}(1-\delta )\right]
\right\} ,  \label{eq:EE_23}
\end{equation}
Marginal product of capital:
\begin{equation}
r_t^K = p_t^H \alpha \frac{y_t^H}{k_{t-1}}a_{t-1}, \label{eq:EE_24}
\end{equation}
Marginal product of labor:
\begin{equation}
w_t = p_t^H (1-\alpha-\gamma) \frac{y_t^H}{h_t}, \label{eq:EE_25}
\end{equation}
Ricardian investment first order condition:
\begin{eqnarray}
\frac{p_t}{q_{t}} &=&\left[ 1-\frac{\gamma_k }{2}\left( \frac{i^R_{t}}{i^R_{t-1}}%
a_{t-1}-\bar{a}\right) ^{2}-\gamma_k \left( \frac{i^R_{t}}{i^R_{t-1}}a_{t-1}-\bar{a%
}\right) \frac{i^R_{t}}{i^R_{t-1}}a_{t-1}\right] u_{t}  \notag \\
&&+\frac{\beta }{a_{t}^{\sigma }}\gamma_k E_{t}\left\{ \frac{v_{t+1}}{v_{t}}%
\frac{\lambda _{t+1}}{\lambda _{t}}\frac{q_{t+1}}{q_{t}}\left( \frac{i^R_{t+1}%
}{i^R_{t}}a_{t}-\bar{a}\right) \left( \frac{i^R_{t+1}}{i^R_{t}}a_{t}\right)
^{2}u_{t+1}\right\} ,  \label{eq:EE_26}
\end{eqnarray}%
Foreign good market clearing:
\begin{equation}
y_{t}^{F}=x_{t}^{F}.  \label{eq:EE_27}
\end{equation}%
Defenition of Ricardian consumption bundle:
\begin{equation}
 \hat{c}_t^R = \left[(1-o_c)^{\frac{1}{\eta_c}}\left(c_t^R-\varsigma \frac{c_{t-1}^R}{a_{t-1}}\right)^{\frac{\eta_c-1}{\eta_c}}+o_c^{\frac{1}{\eta_c}}\left(g_t^c\right)^{\frac{\eta_c-1}{\eta_c}}\right]^{\frac{\eta_c}{\eta_c-1}} 
 \label{eq:EE_28}
\end{equation}
Defenition of hand-to-mouth consumption bundle:
\begin{equation}
\hat{c}_t^{hm} = \left[(1-o_c)^{\frac{1}{\eta_c}}\left(c_t^{hm}-\varsigma \frac{c_{t-1}^{hm}}{a_{t-1}}\right)^{\frac{\eta_c-1}{\eta_c}}+o_c^{\frac{1}{\eta_c}}\left(g_t^c\right)^{\frac{\eta_c-1}{\eta_c}}\right]^{\frac{\eta_c}{\eta_c-1}} 
 \label{eq:EE_29}
\end{equation}
Ricardian wealth effect parameter 1:
\begin{equation}
\Theta^R _{t}=\widetilde{\chi }_{t}^R\left(\left[(1-o_c)^{\frac{1}{\eta_c}}\left(c^R_{t}-\varsigma \frac{c^R_{t-1}}{%
a_{t-1}}\right)^{\frac{\eta_c - 1}{\eta_c}} + o_c^{\frac{1}{\eta_c}}\left(\bar{g}^c\right)^{\frac{\eta_c -1}{\eta_c}}\right]^{\frac{\eta_c}{\eta_c -1}}\right)^{-\sigma}.  \label{eq:EE_30}
\end{equation}%
Ricardian wealth effect parameter 2:
\begin{equation}
\widetilde{\chi }_{t}^R=\widetilde{\chi }_{t-1}^{o^{1-\upsilon} }\left(\left[(1-o_c)^{\frac{1}{\eta_c}}\left(c^R_{t}-\varsigma \frac{c^R_{t-1}}{%
a_{t-1}}\right)^{\frac{\eta_c - 1}{\eta_c}} + o_c^{\frac{1}{\eta_c}}\left(\bar{g}^c\right)^{\frac{\eta_c -1}{\eta_c}}\right]^{\frac{\eta_c}{\eta_c -1}}\right)^{\sigma \upsilon }
\label{eq:EE_31}
\end{equation}
Hand-to-mouth wealth effect parameter 1:
\begin{equation}
\Theta^{hm} _{t}=\widetilde{\chi }_{t}^R\left(\left[(1-o_c)^{\frac{1}{\eta_c}}\left(c^{hm}_{t}-\varsigma \frac{c^{hm}_{t-1}}{%
a_{t-1}}\right)^{\frac{\eta_c - 1}{\eta_c}} + o_c^{\frac{1}{\eta_c}}\left(\bar{g}^c\right)^{\frac{\eta_c -1}{\eta_c}}\right]^{\frac{\eta_c}{\eta_c -1}}\right)^{-\sigma}.  \label{eq:EE_32}
\end{equation}%
Hand-to-mouth wealth effect parameter 2:
\begin{equation}
\widetilde{\chi }_{t}^{hm}=\widetilde{\chi }_{t-1}^{r^{1-\upsilon} }\left(\left[(1-o_c)^{\frac{1}{\eta_c}}\left(c^{hm}_{t}-\varsigma \frac{c^{hm}_{t-1}}{%
a_{t-1}}\right)^{\frac{\eta_c - 1}{\eta_c}} + o_c^{\frac{1}{\eta_c}}\left(\bar{g}^c\right)^{\frac{\eta_c -1}{\eta_c}}\right]^{\frac{\eta_c}{\eta_c -1}}\right)^{\sigma \upsilon }
\label{eq:EE_33}
\end{equation}%
Government budget constraint:
\begin{equation}
p_t g^c_{t} + p_tg^i_{t} +R_{t-1}^*\xi_{t-1}\frac{b^*_{t-1}}{a_{t-1}} + tr_t = \chi p_t^{Co*}y_t^{Co}+b_t^* + \tau_t^n w_t (1-\omega)h_t^R + \tau_t^k (r_t^k-\delta)\frac{k_{t-1}}{a_{t-1}} + \tau_t^c p_t c_t + (1-\omega)T
\label{eq:EE_34}
\end{equation}%
Transfers to Ricardian households:
\begin{equation}
tr_t^R = \frac{1-\omega^G}{1-\omega}tr_t  
\label{eq:EE_35}
\end{equation}
Transfers to hand-to-mouth households:
\begin{equation}
tr_t^{hm} = \frac{\omega^G}{\omega}tr_t  
\label{eq:EE_36}
\end{equation}
Home good producer profits
\begin{equation}
\pi_t = \pi^H_t = p_t^H\gamma y_t^H
\label{eq:EE_37}
\end{equation}%
Government capital evolution:
\begin{equation}
k^g_{t} = (1-\delta_g)\frac{k^g_{t-1}}{a_{t-1}} + g^i_t
\label{eq:EE_38}
\end{equation}%
Consumption tax:
\begin{equation}
\tau_t^c = \tilde{\tau}^c
\label{eq:EE_39}
\end{equation}%
Labor income tax:
\begin{equation}
\tau_t^n = \tilde{\tau}^n
\label{eq:EE_40}
\end{equation}%
Capital income tax:
\begin{equation}
\tau_t^k = \tilde{\tau}^k
\label{eq:EE_41}
\end{equation}%
Government investment spending rule:
\begin{equation}
 g_t^i =\left( g^i_{t-1}\right)^{\rho_{g A}}\left(\bar{g}^i\left(\frac{y_t}{y}\right)^{\alpha_{gi}}\left(\frac{b_{t-1}^{\ast}/y_{t-1}}{\bar{b}}\right)^{-\gamma_{gi}}\right)^{1-\rho_{g i}}\exp(e_t^{gi})  
 \label{eq:EE_42}
\end{equation}
Consumption spending investment rule:
\begin{equation}
 g_t^c =\left( g^c_{t-1}\right)^{\rho_{g c}}\left(\bar{g}^c\left(\frac{y_t}{y}\right)^{\alpha_{gc}}\left(\frac{b_{t-1}^{\ast} /y_{t-1}}{\bar{b}}\right)^{-\gamma_{gc}}\right)^{1-\rho_{g c}}\exp(e_t^{gc})   
 \label{eq:EE_43}
\end{equation}
Transfers spending rule:
\begin{equation}
tr_t =\left( tr_{t-1}\right)^{\rho_{tr}}\left(\bar{tr}\left(\frac{y_t}{y}\right)^{\alpha_{tr}}\left(\frac{b_{t-1}^{\ast}/y_{t-1}}{\bar{b}}\right)^{-\gamma_{tr}}\right)^{1-\rho_{tr}}\exp(e_t^{tr})   
 \label{eq:EE_44}
\end{equation}
Hand-to-mouth budget constraint:
\begin{equation}
(1+\tau_t^c)p_tc_t^{hm} = w_th_t^{hm} + tr_t^{hm}
\label{eq:EE_45}
\end{equation}  
The exogenous processes are%
\begin{equation*}
\log \left( x_{t}/\bar{x}\right) =\rho _{x}\log \left( x_{t-1}/\bar{x}%
\right) +\varepsilon _{t}^{x},\qquad \rho _{x}\in \lbrack 0,1),\qquad \bar{x}%
>0,
\end{equation*}%
for each exogenous sequence $x_t$, where the $\varepsilon _{t}^{x}$ are i.i.d. shocks.
\section{Steady State}
\label{Appendix:steadystate}

We show how to compute the steady state for given values of $h^R$, $h^{hm}$, $p$, $i_r = g^i/i$, $s^{f}=f^*/y$, $s^{g}=pg^c/y$, $s^P=p^{Co\ast }y^{Co}/y$, $\bar{b}=b^{\ast}/y $,$s^{tr} = tr/y$, and $\xi = \bar{\xi}$. The parameters $\beta $, $\bar{\kappa}^R$,$\bar{\kappa}^{hm}$, $o^{\ast }$, $\bar{g}^c$, $\bar{g}^A$, $(\bar{f}^{\ast} + \bar{b}^{\ast})$,$T$, $\bar{tr}$, and $%
\bar{y}^{Co}$ are determined endogenously while the values of the remaining
parameters are exogenous. All exogenous variables are equal to their ``bar" values in steady state.

\noindent From (\ref{eq:EE_7}),%
\begin{equation*}
\beta =\bar{a}^{\sigma }/(\bar{R}^{\ast }\bar{\xi} ).
\end{equation*}%
From (\ref{eq:EE_6})
\begin{equation*}
 h = (1-\omega)h^R + \omega h^{hm}.   
\end{equation*}
From (\ref{eq:EE_26}) and the fact that we calibrate $\bar{u}$ to $1$,%
\begin{equation*}
q=p.
\end{equation*}%
\noindent From (\ref{eq:EE_39}) - (\ref{eq:EE_41})
\begin{align*}
\tau^k &= \bar{\tau}^k\\
\tau^n &= \bar{\tau}^n\\
\tau^c &= \bar{\tau}^c
\end{align*}
From (\ref{eq:EE_23}),%
\begin{equation*}
r^{K}=\frac{q\left( \frac{\bar{a}^{\sigma }}{\beta }-1+\delta \right)-\tau^k\delta}{1-\tau^k}.
\end{equation*}%
Plugging (\ref{eq:EE_9}) and (\ref{eq:EE_10}) in (\ref{eq:EE_8}),
\begin{equation*}
p^H = p\left[\frac{1}{1-o}\left(1-o\left(\frac{1}{p}\right)^{1-\eta}\right)\right]^{\frac{1}{1-\eta}}.
\end{equation*}%
We now perform the following steps to derive $y^H,k,k^g,i,g^i,i^R,$and $k^R.$

\noindent From (\ref{eq:EE_38})
\begin{equation*}
k^g = \frac{g^i}{1-\frac{(1-\delta_g)}{\bar{a}}}.
\end{equation*}%

\noindent From (\ref{eq:EE_20}) and $\bar{u} = 1$,%
\begin{equation*}
k^R=\frac{i^R}{ 1-\frac{1-\delta} {\bar{a}}}.
\end{equation*}%

\noindent From (\ref{eq:EE_21}),
\begin{equation*}
k = k^R(1-\omega).
\end{equation*}\noindent From (\ref{eq:EE_22}),
\begin{equation*}
i^R = i/(1-\omega)
\end{equation*}
Combining these results yields 
\begin{equation*}
 k_r \equiv \frac{k^g}{k} = i_r\frac{1-\frac{1-\delta}{\bar{a}}}{1-\frac{1-\delta_g}{\bar{a}}}.   
\end{equation*}
This result along with equations (\ref{eq:EE_12}) and (\ref{eq:EE_24}) form a system in $y^H$, $k$, and $k^g$ which gives
\begin{equation*}
y^H = \left(\bar{z}\left(k_r\frac{ p^H\alpha}{r^k}\right)^\gamma\left(\frac{p^H\alpha}{r^K}\right)^\alpha(\bar{a}h)^{1-\alpha-\gamma}\right)^{\frac{1}{1-\alpha-\gamma}}
\end{equation*}%
\begin{equation*}
k = p^H \alpha \frac{y^H}{r^K}\bar{a},
\end{equation*}%
\begin{equation*}
k^g = k_rk.
\end{equation*}
Then substituting in these results yields 
\begin{equation*}
 g^i = \left(1-\frac{1-\delta_g}{\bar{a}}\right)k^g 
\end{equation*}
\begin{equation*}
k^R = k/(1-\omega),    
\end{equation*}
\begin{equation*}
 i^R = \left(1-\frac{1-\delta}{a}\right)k^R,
\end{equation*}
\begin{equation*}
i = (1-\omega)i^R.    
\end{equation*}
From (\ref{eq:EE_25}),
\begin{equation*}
w = p^H(1-\alpha-\gamma)\frac{y^H}{h}.
\end{equation*}%
Equations (\ref{eq:EE_15}) and the definition of the share of commodity output in GDP, $s^P=(p^{Co\ast}y^{Co})/y$, form a system in $y$ and $\bar{y}^{Co}$, which gives
\begin{equation*}
y = \frac{p^H y^H}{1-s^{Co}}
\end{equation*}%
\begin{equation*}
\bar{y}^{Co} = \frac{s^{Co}y}{p^{Co\ast}}.
\end{equation*}%

\noindent From $s^{g}=pg^c/y$,
\begin{equation*}
 g^c=s^{g}y/p.
\end{equation*}%

\noindent From (\ref{eq:EE_42}) and the assumption $\bar{b}=b^*/y$, 
\begin{equation*}
 \bar{g}^i = g^i.   
\end{equation*}
From (\ref{eq:EE_43}) and the assumption $\bar{b}=b^{\ast}/y$,
\begin{equation*}
 \bar{g}^c = g^c.   
\end{equation*}
From $\bar{b} = b^*/y$
\begin{equation*}
b^* = \bar{b}y
\end{equation*}
From $\bar{f} = f^*/y$,
$$f^* = \bar{f}y.$$
From (\ref{eq:EE_17}),
\begin{equation*}
 f^{o*}  = \frac{f^*}{1-\omega}  
\end{equation*}
From (\ref{eq:EE_18}),
\begin{equation*}
tb = (1-\chi)*p^{Co\star}y^{Co} - \left(1-\frac{\bar{R}^*\bar{\xi}}{\bar{a}}\right)(f^*+b^*) .
\end{equation*}
From (\ref{eq:EE_19}) and $\xi = \bar{\xi}$,
\begin{equation*}
(\bar{f}^{\ast}+\bar{b}^{\ast}) = \frac{f^{\ast}+b^{\ast}}{y}.
\end{equation*}
From the definition of GDP as the sum of domestic absorption and the trade balance (\ref{eq:EE_15}),
\begin{equation*}
c=\frac{y-tb}{p}-i-g^c-g^i.
\end{equation*}
From (\ref{eq:EE_14}),
\begin{equation*}
y^C = c + i + g^c + g^i
\end{equation*}%
From (\ref{eq:EE_9}),
\begin{equation*}
x^F = o\left(\frac{1}{p}\right)^{-\eta}y^C.
\end{equation*}%
From (\ref{eq:EE_10}),
\begin{equation*}
x^H = (1-o)\left(\frac{p^H}{p}\right)^{-\eta}y^C.
\end{equation*}%
From (\ref{eq:EE_27}),
\begin{equation*}
y^F=x^F.
\end{equation*}%
From (\ref{eq:EE_16}),
\begin{equation*}
x^{H\ast}= \frac{tb-\bar{p}^{Co\ast }y^{Co}+y^F}{p^H}.
\end{equation*}
From (\ref{eq:EE_11}),
\begin{equation*}
o^{\ast}=\frac{x^{H\ast }}{\left(p^H\right)^{-\eta \ast }\bar{y}^{\ast}}.
\end{equation*}
From (\ref{eq:EE_37}),
\begin{equation*}
\pi = p^H\gamma Y^H.
\end{equation*}
From the defenition of $s^{tr}$,
\begin{equation*}
tr = s^{tr}y
\end{equation*}
From (\ref{eq:EE_34}),
\begin{equation*}
T = \frac{\chi p^{Co*}y^{Co} - R^*\bar{\xi}\frac{b^*}{\bar{a}}  +\tau^n W(1-\omega)h^0 + \tau^k(r^k-\delta)\frac{k}{\bar{a}}+\tau^c pc + \tau^{\pi}\pi - pg^c - pg^i + b^* - tr}{1-\omega}.
\end{equation*}
From (\ref{eq:EE_44}),
\begin{equation*}
 \bar{tr} = tr.   
\end{equation*}
From (\ref{eq:EE_35})-(\ref{eq:EE_36}),
\begin{equation*}
 tr^R = \frac{1-\omega^G}{1-\omega}tr  
\end{equation*}
\begin{equation*}
tr^{hm} = \frac{\omega^G}{\omega}tr.
\end{equation*}
From (\ref{eq:EE_45})
\begin{equation*}
c^{hm} = \frac{wh^{hm} + tr^{hm}}{(1+\tau^c)p}.    
\end{equation*}
From (\ref{eq:EE_3}),
\begin{equation*}
 c^R = \frac{c-\omega c^{hm}}{1-\omega}.   
\end{equation*}
From (\ref{eq:EE_28}),
\begin{equation*}
 \hat{c}^R = \left[(1-o_c)^{\frac{1}{\eta_c}}\left(c^R-\varsigma \frac{c^R}{\bar{a}}\right)^{\frac{\eta_c-1}{\eta_c}}+o_c^{\frac{1}{\eta_c}}\left(g^c\right)^{\frac{\eta_c-1}{\eta_c}}\right]^{\frac{\eta_c}{\eta_c-1}}. 
\end{equation*}
From (\ref{eq:EE_29}),
\begin{equation*}
\hat{c}^{hm} = \left[(1-o_c)^{\frac{1}{\eta_c}}\left(c^{hm}-\varsigma \frac{c^{hm}}{\bar{a}}\right)^{\frac{\eta_c-1}{\eta_c}}+o_c^{\frac{1}{\eta_c}}\left(g^c\right)^{\frac{\eta_c-1}{\eta_c}}\right]^{\frac{\eta_c}{\eta_c-1}}.
\end{equation*}
From (\ref{eq:EE_31}),%
\begin{equation*}
\widetilde{\chi}^R=\hat{c}^{o^\sigma}.
\end{equation*}%
From (\ref{eq:EE_33}),%
\begin{equation*}
\widetilde{\chi}^{hm}=\hat{c}^{r^\sigma}.
\end{equation*}%
From (\ref{eq:EE_30}),%
\begin{equation*}
\Theta^R =1.
\end{equation*}%
From (\ref{eq:EE_32}),%
\begin{equation*}
\Theta^{hm} =1.
\end{equation*}%
From (\ref{eq:EE_1}),
\begin{equation*}
\lambda^R=\frac{\hat{c}^{o^{-\sigma}}}{(1+\tau^c)p}\left(\frac{(1-o_c)\hat{c}^R}{c^R - \varsigma\frac{c^R}{\bar{a}} }\right)^{\frac{1}{\eta_c}}.
\end{equation*}
From (\ref{eq:EE_2}),
\begin{equation*}
\lambda^{hm}=\frac{\hat{c}^{r^{-\sigma}}}{(1+\tau^c)p}\left(\frac{(1-o_c)\hat{c}^{hm}}{c^{hm} - \varsigma\frac{c^{hm}}{\bar{a}} }\right)^{\frac{1}{\eta_c}}.
\end{equation*}
From (\ref{eq:EE_4})
\begin{equation*}
\bar{\kappa}^R=\frac{\lambda^R (1-\tilde{\tau}^n)w}{h^{o^\phi}}. 
\end{equation*}
From (\ref{eq:EE_5})
\begin{equation*}
\bar{\kappa}^{hm}=\frac{\lambda^{hm} w}{h^{r^\phi}}. 
\end{equation*}

\section{Calibrated Parameters}\label{app:calibratedparameters}
Following \citet{MedinaSoto2016}, we use logarithmic preferences, setting the inverse elasticity of intertemporal substitution, $\sigma=1$. We follow \citet{coenen2013gauging} and set the share of private consumption in the aggregate consumption bundle, $1-o_c$, equal to 0.75. The depreciation rates of private and public capital, $\delta$ and $\delta_g$, are set to 6\% annually, following \citet{MedinaSoto2016} for the former and \citet{coenen2013gauging} for the latter. The parameter characterizing home bias in domestic demand, $o$, is set to 0.27, which is the average contribution of imports over domestic demand in the 2000\textendash2019 period. In our Cobb-Douglas specification for the production of home goods, the output elasticity of labor, $1-\alpha-\gamma$, is 0.66, following \citet{MedinaSoto2016}. The share of public capital is set to 0.034. This value corresponds to 10\% of the share of total capital (private plus public, $\alpha+\gamma=0.34$), and is motivated by \citet{coenen2013gauging}, who use a share of 0.1 for public capital in their CES aggregator of private and public capital. 

We set the price elasticity of foreign demand for home goods, $\eta^*$, to 0.25 \citep[see the estimates in][]{crespo2004estimacion, pelaez2020elasticidades}. The government's share of the income generated in the commodity sector, $\chi$, is set to 0.74, considering that the average share of total oil output produced by the public oil company during 2000\textendash 2019 is 60\%, and a general tax on foreign firms is 35\%.\footnote{Note that $\chi= c + (1-c)*t$, where $c$ is the share of production owned by the government and $t$ is the average tax rate on foreign oil producers.} The share of hand-to-mouth households, $\omega$, is set to 0.7, following estimates by the Central Bank of Ecuador that suggest that about 30\% of the population has an outstanding loan in the financial system \citep[see][]{BCE2023EIF}. We assume that all government transfers go to the hand-to-mouth households, setting $\omega_G=1$. Finally, normalizing the time endowment of households to unity and using the historical average of hours worked per week adjusted by the average employment rate, we set $h^R=h^{hm}=0.3$.

Regarding the targeted first moments, most values are drawn from the 2022-2027 projections in the July 2022 IMF Staff Report for Ecuador \citep[see][]{oner2022imf}. This is the case for the ratios of government consumption and transfers to GDP, $s^g$ and $s^{tr}$, the ratios of private and public external debt to GDP, $\bar{f}$ and $\bar{b}$, and the share of oil production in GDP, $s^P$. Other targeted moments are sample averages or drawn from other sources. This is the case for the ratio of public to total investment, $i_r$, the annualized quarterly growth rate of labor productivity, $\bar{a}$, the annualized quarterly gross real 3-month T-bill rate, $R^*$, the country risk premium, $\bar{\xi}$, and the consumption, corporate, and labor income tax rates, $\bar{\tau}_c$, $\bar{\tau}_k$, $\bar{\tau}_n$.

We also calibrate the parameters associated with the exogenous processes for which we have a data counterpart\textemdash oil production, the international interest rate, world GDP, and the price of oil. In all cases, we remove trends from these variables using the \citet{hodrick1997filter} filter and estimate univariate AR(1) processes for the same sample period we use to estimate the rest of the parameters of the model.

%\begin{landscape}
%\small
\begin{table}
\begin{footnotesize}
\begin{center}
    \caption{Calibrated Parameters and Targeted Steady State Values}\label{table:steadystatecalibration}
    \begin{tabular}{c l c l}
        \hline
        Parameter & Description & Value & Source \\ \hline
        $\sigma$ & Inverse elast. intert. subst. & 1 & \citet{MedinaSoto2016} \\
        ($1-o_c$) & Share of private consumption in $\hat{C}$ & 0.75 & \citet{coenen2013gauging} \\
        $\delta$ & Annualized depreciation rate ($K$) & 6\% & \citet{MedinaSoto2016} \\
        $\delta_g$ & Annualized depreciation rate ($K^g$) & 6\% & \citet{coenen2013gauging} \\
        $o$ & Home bias in domestic demand & 0.27 & Avg. share of imports (2000\textendash2019) \\
        $1-\alpha-\gamma$ & Output elasticity of labor & 0.66 & \citet{MedinaSoto2016} \\
        $\gamma$ & Output elasticity of public capital & 0.034 & \citet{coenen2013gauging} \\
        $\eta^*$ & Price elasticity of foreign demand & 0.25 & \citet{pelaez2020elasticidades} \\
        $\chi$ & Government share in commodity sector & 0.74 & Average (2000\textendash 2019) \\
        $\omega$ & Share of hand-to-mouth households & 0.7 & Central Bank of Ecuador estimates \\
        $h^R$, $h^{hm}$ & Hours worked & 0.3 & Assumption and INEC statistics\\
        \hline
        $s^g$ & Government consumption to GDP ratio & 0.11 & \citet{oner2022imf} \\
        $s^{tr}$ & Government transfers to GDP ratio & 0.05 & \citet{oner2022imf} \\
        $\bar{f}$ & Private external debt to quarterly GDP ratio & 0.26 & \citet{oner2022imf} \\
        $\bar{b}$ & Total external debt to quarterly GDP ratio & 1.42 & \citet{oner2022imf} \\
        $s^P$ & Oil production to total GDP ratio & 0.08 & \citet{oner2022imf} \\
        $i_r$ & Public to total investment ratio & 0.16 & OECD average (2007-2020) \\
        $\bar{a}$ & Annualized balanced growth path & 2\% & \citet{gonzalez2016como} \\
        $R^*$ & Quarterly gross real foreign int. rate & 1.001 & \citet{cbo2021projections} \\ 
        $\bar{\xi}$ & Quarterly gross country premium & 1.013 & Latin America avg. (1997-2020) \\ %Quarterly gross (real) sovereign bond spread
        $\bar{\tau}_c$ & Consumption tax rate & 12\% & Current rate \\
        $\bar{\tau}_k$ & Corporate tax rate & 25\% & Current rate \\
        $\bar{\tau}_n$ & Labor income tax rate & 4\% & Avg. labor income tax rate \\
        \hline
    \end{tabular}
\end{center}
\end{footnotesize}
\end{table}

\section{Estimated Parameters}\label{app:estimatedparameters}

The rest of the parameters of the model are estimated using Bayesian techniques, solving the model with a linear approximation around the non-stochastic steady state. For the observable variables, we use quarterly information from Ecuador for the period 2004\textendash 2019. Ecuador adopted the U.S. dollar as its official currency in January 2000. Our sample begins in 2004 because inflation took about 4 years to stabilize after a banking crisis in 1999 and the introduction of the U.S. dollar. We use the following 13 observable variables: Non-oil GDP, private aggregate consumption, aggregate investment, government consumption, government investment, government transfers, the trade balance, hours worked, the country risk premium as measured by the Emerging Markets Bond Index (EMBI) for Ecuador, the West Texas Intermediate (WTI) oil price, oil production, the foreign interest rate, and foreign GDP. Because the foreign good is the numeraire in the model, we deflate all the variables, which are expressed in current U.S. dollars, using the U.S. GDP deflator. In addition, we express all the flow variables as quarter-over-quarter growth rates, except hours worked and foreign output, which are expressed as percent deviations from a \citet{hodrick1997filter} filtered trend. We do the same for the foreign interest rate, the (log real) oil price, and oil production. As we previously mentioned, we calibrate the parameters of the exogenous processes for which we have a data counterpart, but we also include these variables as observables in the Bayesian estimation.\footnote{While the parameters of these exogenous processes were calibrated, including these variables in the data set is informative for the inference of the innovations associated with the other exogenous processes.} Our estimation strategy also includes i.i.d. measurement errors for all observables. The variance of the measurement errors are set to 10\% of the variance of the corresponding observables. Appendix \ref{appendix:datadetails} contains the sources of the data and offers details on the construction of the observable variables.

Table \ref{table:priorposterior} shows the prior and posterior distributions of the estimated parameters. The prior distributions are fairly loose. For the parameters of the fiscal expenditure rules ($\rho_x, \alpha_x, \gamma_x$ for $x = gc, gi, TR$), we center the prior distributions using coefficients obtained from auxiliary OLS regressions of the corresponding rules on measures of spending, output, and public debt gaps. These regressions are estimated outside the structural model and serve only to anchor the prior means in empirically plausible regions of the parameter space. The priors are specified to be sufficiently diffuse so that the posterior estimates are primarily informed by the likelihood of the full model. The prior distributions of the persistence coefficients of the structural shocks ($\rho_x$ for $x=\upsilon, \kappa, u, z, a, \zeta$) are centered at 0.75, following the DSGE model estimated by \citet{guerra2021search} for Chile. We follow the same paper to inform the prior mean of the elasticity of the country risk premium to the aggregate debt-to-GDP ratio, $\psi$, which is centered at 0.005. The coefficient of investment adjustment costs, $\gamma_k$, is centered at the value obtained by \citet{smets2007shocks}, close to 6. The rest of the parameters, including the variances of the shocks, were chosen to loosely match the sample standard deviation of the observed variables with their model counterparts under the prior means of the other coefficients, but the standard deviations of the prior distributions are wide enough to allow the data to inform the estimated posterior distributions.

\begin{table}
\begin{footnotesize}
    \centering
    \caption{Estimated Parameters}
    \label{table:priorposterior}
    \begin{tabular}{c l | l c c}
    \hline \hline
    Parameter & \multicolumn{1}{c}{Description} & \multicolumn{1}{c}{Prior Distribution} & Post. Mean & Post. SD \\
    \hline 
    $\phi$ & Inv. Frisch elast. of $h$ & $N^{(0,\infty)}\left(0.1,2\right)$ & 3.42 & 1.00  \\
    $\varsigma$ & Habit formation & $N^{(0,1)}\left(0.9,0.2\right)$ & 0.74 & 0.09 \\
    $\eta_c$ & Elast. of subst. $(C,G^c)$ & $N^{(0,\infty)}\left(1.5,2\right)$ & 0.79 & 0.71 \\
    $\nu$ & Wealth effect & $N^{(0,1)}\left(0.5,0.5\right)$ & 0.70 & 0.20 \\
    $\eta$ & Elast. subst. home-foreign & $N^{(0,\infty)}\left(3,2\right)$ & 6.08 & 1.06 \\
    $\gamma_k$ & Inv. adj. cost & $N^{(0,\infty)}\left(6,10\right)$ & 8.85 & 4.44  \\
    100$\psi$ & Country prem. debt elast. & $N^{(0,\infty)}\left(0.5,10\right)$ & 2.21 & 1.40 \\
    \hline
    $\rho_{gc}$ & Gov. cons. persistence & $N^{(0,1)}\left(0.70,1\right)$ & 0.88 & 0.07 \\
    $\rho_{gi}$ & Gov. inv. persistence & $N^{(0,1)}\left(0.48,1\right)$ & 0.77 & 0.08 \\
    $\rho_{TR}$ & Transfers persistence & $N^{(0,1)}\left(0.37,1\right)$ & 0.47 & 0.12 \\
    $\alpha_{gc}$ & Gov. cons. cyclicality & $N\left(0.13,2\right)$ & 0.59 & 0.92 \\
    $\alpha_{gi}$ & Gov. inv. cyclicality & $N\left(1.77,2\right)$ & 2.58 & 1.57 \\
    $\alpha_{TR}$ & Transfers cyclicality & $N\left(3.52,2\right)$ & 5.46 & 0.96 \\
    $\gamma_{gc}$ & Gov. cons. debt stabilization & $N^{(-\infty,0)}\left(-0.01,2\right)$ & -0.50 & 0.41 \\
    $\gamma_{gi}$ & Gov. inv. debt stabilization & $N^{(-\infty,0)}\left(-0.13,2\right)$ & -1.82 & 0.67 \\
    $\gamma_{TR}$ & Transfers debt stabilization & $N^{(-\infty,0)}\left(-0.35,2\right)$ & -0.26  & 0.21 \\
    \hline
    $\rho_{\upsilon}$ & AC cons. pref. sh. & Beta$\left(0.75,0.15\right)$ & 0.41 & 0.12 \\
    $\rho_{\kappa}$ & AC labor pref. sh. & Beta$\left(0.75,0.15\right)$ & 0.74 & 0.10 \\
    $\rho_{u}$ & AC invest. sh. & Beta$\left(0.75,0.15\right)$ & 0.35 & 0.11 \\
    $\rho_{z}$ & AC stationary tech. sh. & Beta$\left(0.75,0.15\right)$ & 0.86 & 0.09 \\
    $\rho_{a}$ & AC non-stationary tech. sh. & Beta$\left(0.75,0.15\right)$ & 0.79 & 0.13 \\
    $\rho_{\zeta}$ & AC country prem. sh. & Beta$\left(0.75,0.15\right)$ & 0.65 & 0.06 \\
    $\sigma_{e^{gc}}$ & SD gov. cons. sh. & IG$\left(0.03,\infty\right)$ & 0.03 & 0.00 \\
    $\sigma_{e^{gi}}$ & SD gov. inv. sh. & IG$\left(0.15,\infty\right)$ & 0.17 & 0.02 \\
    $\sigma_{e^{TR}}$ & SD transfers sh. & IG$\left(0.10,\infty\right)$ & 0.11 & 0.01 \\
    $100\sigma_{\upsilon}$ & SD cons. pref. sh. & IG$\left(1,\infty\right)$ & 10.23 & 2.63 \\
    $100\sigma_{\kappa}$ & SD labor pref. sh. & IG$\left(4,\infty\right)$ & 7.56 & 1.91 \\
    $100\sigma_{u}$ & SD invest. sh. & IG$\left(1,\infty\right)$ & 27.24 & 13.09 \\
    $100\sigma_{z}$ & SD stationary tech. sh. & IG$\left(0.1,\infty\right)$ & 2.08 & 0.30 \\
    $100\sigma_{a}$ & SD non-stationary tech. sh. & IG$\left(0.1,\infty\right)$ & 0.62 & 0.39 \\
    $100\sigma_{\zeta}$ & SD country prem. sh. & IG$\left(0.9,\infty\right)$ & 0.89 & 0.08 \\
    \hline
    \end{tabular}
    \caption*{\tiny{\emph{Note:} $N^{(l,u)}\left(a,b\right)$ denotes a normal distribution with mean $a$ and standard deviation $b$ that was restricted to take values between $l$ and $u$ during mode finding, but not during posterior distribution sampling. Beta$\left(a,b\right)$ denotes a beta distribution with mean $a$ and standard deviation $b$. $IG(a,b)$ denotes an inverse gamma distribution for the variance of the shocks where $a$ is the mean of the standard deviation and $b$ is its variance. Results are based on two Metropolis-Hastings chains of 500,000 draws each from the posterior distribution, after dropping 500,000 draws for each chain, using an acceptance rate of 0.3, and starting the chains around the posterior mode obtained with Dynare's Monte-Carlo\textendash based optimization routine.}}
\end{footnotesize}
\end{table}
%%%%%%%%%%%%%%%%%%%%%%%%%%%%%%%%%%%%%%%%%%%%%%%%%%%%%%

We now discuss the results of the Bayesian estimation. The posterior mean estimates of the parameters associated with the fiscal expenditure rules suggest all three components (government consumption, investment, and transfers) are procyclical, because the parameters $\alpha_{gc}$, $\alpha_{gi}$ and $\alpha_{TR}$ are all greater than zero. Fiscal transfers are the most procyclical component, whereas government investment is the least procyclical. Regarding the debt-stabilization parameters of the fiscal expenditure rules, we find that government investment is the most active expenditure instrument in debt stabilization. The negative values of these parameters indicate that when the public debt-to-GDP ratio deviates from target, expenditure moves in the opposite direction.\footnote{Note that the prior distributions of the debt-stabilization parameters are truncated to avoid de-stabilizing values.}

The second set of results concerns preference parameters. The estimated value of $\phi$, at 3.4, suggests a low elasticity of labor supply. The wealth effect on labor supply, $\nu$, at 0.7, is estimated to be closer to CRRA formulations (the case when $\nu=1$), than formulations that shut down the wealth effect (the case when $\nu=0$). Additionally, the habit persistence parameter, $\varsigma$, at 0.7, is estimated in the vicinity of the results in the literature. The parameter that governs the elasticity of substitution between private and government-provided goods, $\eta_c$, at 0.8, is estimated to be relatively low, pointing to complementarity between these two types of goods. Finally, the consumption preference shock, $\upsilon$ is estimated to be less persistent than the labor preference shock, $\kappa$, although it is also more volatile. 

A third set of parameters is related to production and investment technologies. For instance, the estimation points to a high degree of substitutability between home and foreign goods, as indicated by the posterior mean of the parameter $\eta$. The investment adjustment cost parameter, $\gamma_k$, is in the upper bound of the results in the literature for developing economies \citep[see, for example][]{sheen2016assessing, fornero2018learning, guerra2021search}, reflecting the higher investment adjustment costs faced by countries dependent on commodities with volatile prices, as argued by \citet{devlin2004managing}. Among technology shocks, the investment efficiency shock, $u_t$, is estimated to have the lowest persistence and highest variance. The shock to the non-stationary technological index, $a_t$, is estimated to be highly persistent for a shock that generates permanent effects on the level of variables in the economy. As we explain below, this is a reasonable inference given the performance of the Ecuadorian economy in the sample period.

The last set of coefficients is related to the country risk premium, $\xi$. The posterior mean of parameter $\psi$ suggests that a 1 percentage point increase in the total foreign debt-to-GDP ratio is estimated to increase the premium by about 23 basis points.\footnote{To gain intuition on this interpretation, note that linearizing the expression for the country premium, \eqref{eq:country_premium}, around $\bar{f}^{\ast} + \bar{b}^{\ast}$, and ignoring the idiosyncratic risk shock, we get $\xi _{t}=\bar{\xi} + \frac{\psi}{\bar{f}^{\ast} + \bar{b}^{\ast}}((F^{\ast}_t +B_{t}^{\ast})/Y_{t}-(\bar{f}^{\ast} + \bar{b}^{\ast}))$. Hence, the effect of a marginal increase in the quarterly debt-to-GDP ratio implies an increase in the quarterly risk premium equal to $\frac{\psi}{\bar{f}^{\ast} + \bar{b}^{\ast}}$. Consequently, a  1 pp increase in the annualized debt-to-GDP ratio is equivalent to an increase of 0.04 in $(F^{\ast}_t +B_{t}^{\ast})/Y_{t}$, which is the quarterly counterpart. Then, the effect on the quarterly country premium is $0.04 \times \frac{\psi}{\bar{f}^{\ast} + \bar{b}^{\ast}}$, where $\bar{f}^{\ast} + \bar{b}^{\ast}=1.55$. As a consequence, the effect is $0.04\times 0.0221/1.55 = 0.00057$. Annualizing this effect results in $1.00057^4-1=0.0023$, i.e., 23 basis points. Interestingly, a linear regression of the EMBI spread against the total foreign debt-to-GDP ratio results in a coefficient equal to 19.}  The idiosyncratic risk premium shock, $\zeta$, displays moderate levels of persistence and volatility, relative to the other shocks.

%%%%%%%%%%%%%%%%%%%%%%%%%%%%%%%%%%%%%%%%%%%%%%%%%%%%%%

\section{Data Details}
\label{appendix:datadetails}
This section describes the data used, their sources, and the transformations we performed in order to estimate the model. The variables are as follows:
\begin{itemize}
    \item Non oil GDP in current dollars: From the Central Bank of Ecuador's (CBE) national accounts bulletin. We deflate the data with the U.S. GDP price deflator and obtain quarter-over-quarter percent growth rates.
    \item Household consumption in current dollars: From the CBE's national accounts bulletin. We deflate the data with the U.S. GDP price deflator and obtain quarter-over-quarter percent growth rates.
    \item Total investment in current dollars: From the CBE's national accounts bulletin. We deflate the data with the U.S. GDP price deflator and obtain quarter-over-quarter percent growth rates.
    \item Oil production in current dollars: From the CBE's national accounts bulletin. We divide the oil GDP by the quarterly average oil price and apply the Hodrick and Prescott filter with $\lambda=1600$ to obtain the cycle of that series in logs.
    \item Hours worked in current dollars: From the National Institute of Statistics and Censuses (NISC). We adjust the quarterly average weekly hours by the employment rate and divide by 120 hours, which is the total number of hours a week a person could work Monday through Friday. We then apply the Hodrick and Prescott filter with $\lambda=1600$ to obtain the cycle of that series in logs.
    \item Government consumption spending in current dollars: From the CBE's national accounts bulletin. We deflate the data with the U.S. GDP price deflator and obtain quarter-over-quarter percent growth rates.
    \item Government investment spending in current dollars: From the CBE's monthly statistical information. We aggregate the monthly data to quarterly and then deflate the data with the U.S. GDP price deflator and obtain quarter-over-quarter percent growth rates.
    \item Government transfers in current dollars: From the CBE's monthly statistical information. We aggregate the monthly data to quarterly and then deflate the data with the U.S. GDP price deflator and obtain quarter-over-quarter percent growth rates.
    % \item Government external debt in current dollars: From the CBE's monthly statistical information. We aggregate the monthly data to quarterly, obtain the first difference and divide by nominal GDP.
    % \item Private external debt in current dollars: From the CBE's monthly statistical information. We aggregate the monthly data to quarterly, obtain the first difference and divide by nominal GDP.
    \item Trade balance in current dollars: From the CBE's national accounts bulletin. We divide by nominal GDP.
    \item 3-month Treasury bill rate: From the Federal Reserve Economic Database (FRED). We aggregate the monthly data to quarterly and obtain the equivalent quarterly gross rate. We then subtract the quarter-over-quarter U.S. inflation gross rate from the U.S. GDP price deflator.
    \item World GDP: From the OECD Statistics. We take the quarter-over-quarter growth rate of the G20 real GDP, compound the rate to obtain a level series, and then apply the Hodrick and Prescott filter with $\lambda=1600$ to obtain the cycle of that series in logs.
    \item Oil price in current dollars: From the CBE's monthly statistical information. We aggregate the monthly data to quarterly, deflate the data with the U.S. GDP price deflator, and then apply the Hodrick and Prescott filter with $\lambda=1600$ to obtain the cycle of that series in logs.
    \item EMBI: From the CBE's monthly statistical information. We aggregate the monthly data to quarterly, divide by 10,000 and add one to make it gross, and then obtain the equivalent quarterly gross rate.
    \item U.S. GDP price deflator: From the FRED database.
\end{itemize}
We demean all the series.

\section{Model-implied Moments against the Data}\label{app:moments}
The model does a good job of matching the co-movement of most of these variables with GDP, both qualitatively and quantitatively. In terms of standard deviations, the model does a good job of matching the relative volatility of most variables. For example, investment is more volatile than consumption, which in turn is more volatile than output; government investment is the most volatile component of public expenditure, followed by transfers and government consumption. Quantitatively, however, the model overstates the volatility of consumption and investment. Finally, the model fits the persistence of the data very well, because the first-order autocorrelation coefficients are quantitatively close to their data counterparts.

 \begin{table}
\begin{footnotesize}
     \centering 
     \caption{Co-Movement, Volatility and Persistence}  \label{table:data_model_moments}
     \begin{tabular}{l  d{1.2} d{1.2}  d{1.2} d{1.2}  d{1.2} d{1.2} }
     \hline \hline
     & \multicolumn{2}{c}{Corr. with GDP} & \multicolumn{2}{c}{SD} & \multicolumn{2}{c}{1st order AC} \\
     & \multicolumn{1}{c}{Data} & \multicolumn{1}{c}{Model} & \multicolumn{1}{c}{Data} & \multicolumn{1}{c}{Model} & \multicolumn{1}{c}{Data} & \multicolumn{1}{c}{Model} \\  
	\hline
	 Non-Oil GDP Growth	            & 1	    & 1	    & 1.65    & 1.54    & 0.46  & 0.41 \\
	 Private Consumption Growth	    & 0.76  & 0.63  & 1.90    & 2.39    & 0.21  & 0.26 \\
	 Total Investment Growth	    & 0.64  & 0.51  & 3.45    & 5.18    & 0.53  & 0.27 \\
	 Trade Balance to GDP	        & 0.12  & 0.02  & 1.97    & 3.97    & 0.67  & 0.77 \\
	 Hours Worked	                & -0.03 & -0.15 & 1.97    & 0.81    & 0.36  & 0.74 \\
	 Government Consumption Growth	& 0.67  & 0.29  & 3.30    & 3.66    & 0.15  & 0.15 \\
	 Government Investment Growth	& 0.33  & 0.30  & 19.39   & 19.42   & -0.08 & -0.01 \\
	 Government Transfers Growth	& 0.49  & 0.36  & 15.77   & 15.91   & 0.26  & -0.01 \\
	 Country Premium	            & -0.14 & -0.16 & 1.30    & 1.26    & 0.73  & 0.70 \\
\hline \hline
 \end{tabular}
 \caption*{\tiny{\emph{Note}: GDP is non-oil real GDP, SD denotes standard deviation, AC denotes autocorrelation.}}
\end{footnotesize}
\end{table}

We comment on the role of commodity, risk, and fiscal policy shocks, which are key aspects of our analysis. Commodity and risk premium shocks are not important drivers of fluctuations in output growth, but are important drivers of the trade balance-to-GDP ratio due to their direct influence on oil exports and the external cost of financing. Commodity shocks drive 25\% of the fluctuations in government transfers, most likely due to the substantial gasoline subsidies offered to households in Ecuador. We now turn to fiscal policy shocks. These shocks explain a significant fraction of fluctuations in private consumption, likely due to the complementarity with private consumption goods in households' consumption basket. Fiscal policy shocks are also important drivers of investment growth, due to the role of public capital in the demand for private capital of the home goods-producing firm. Finally, as is usually the case in real business cycle models, the most important drivers of aggregate fluctuations are technology shocks, especially for output and investment growth, and preference shocks, especially or consumption growth and hours worked.

\begin{table}[]
\begin{footnotesize}
    \centering
    \caption{Variance Decomposition}
    \label{table:vardecomp}
    \begin{tabular}{l  d{2.1} d{2.1} d{2.1} d{2.1} d{2.1} d{2.1}}
    \hline \hline
                                    & \multicolumn{1}{c}{Commodity} & \multicolumn{1}{c}{Risk} & \multicolumn{1}{c}{FP} & \multicolumn{1}{c}{Technology} & \multicolumn{1}{c}{Foreign} & \multicolumn{1}{c}{Preferences} \\ %Country risk
    \hline
    Non-Oil GDP Growth              & 3.1  & 1.4  & 6.3  & 73.8 & 0.6 & 14.8 \\
    Private Consumption Growth      & 10.7 & 5.0  & 14.4 & 19.7 & 0.8 & 49.3 \\
    Total Investment Growth         & 3.7  & 2.8  & 27.5 & 63.4 & 0.4 & 2.2 \\
    Trade Balance-to-GDP Ratio      & 33.8 & 11.8 & 6.4  & 31.1 & 1.7 & 15.3 \\
    Hours Worked                    & 4.4  & 1.2  & 2.3  & 9.8  & 0.2 & 82.2 \\
    Government Consumption Growth   & 5.9  & 1.1  & 70.9 & 18.0 & 0.2 & 4.0 \\
    Government Investment Growth    & 4.9  & 0.8  & 91.5 & 2.3  & 0.1 & 0.4 \\
    Government Transfers Growth     & 25.7 & 0.4  & 66.8 & 5.5  & 0.1 & 1.5 \\
    Country Premium	                & 2.8  & 92.7 & 0.3  & 3.2  & 0.0 & 0.3 \\
    \hline \hline
    \end{tabular}
    \caption*{\tiny \emph{Note}: Fraction of the unconditional theoretical variances at the posterior mean (in percent) explained by the shocks, excluding measurement error. Commodity shocks include oil price ($p^{P*}$) and oil output ($Y^{P}$) shocks. The risk shock affects the country premium ($\zeta$). Fiscal policy (FP) shocks include those to government consumption, government investment, and transfers ($e^{gc}$, $e^{gi}$, and $e^{TR}$). Technology shocks include transitory ($z$) and permanent ($a$) productivity shocks, and investment efficiency ($u$) shocks. Foreign shocks include shocks to the foreign interest rate ($R^*$) and foreign demand ($y^*$). Preference shocks include consumption ($v$) and labor supply ($\kappa$) shocks. Columns may not add to their exact percentages due to rounding.}
\end{footnotesize}
\end{table}

\section{Shock Decomposition for Observables}\label{Appendix:shockdecomp}
This appendix shows the historical shock decomposition of the observable variables in figures \ref{fig:shockdecompnonoil}-\ref{fig:shockdecompforeigngdp} through the lens of the model.
\begin{figure}
    \centering
     \caption{Historical Shock Decomposition for Non-Oil GDP Growth Rate} \label{fig:shockdecompnonoil}
    \includegraphics [scale=.4]{Figures/gam_YNCo_obs}
    \smallskip
    \captionsetup{width=.8\linewidth}
    \caption*{{\tiny The black line is the observed variable. The y-axis is percentage point deviation from steady state and the x-axis is quarters.  }}
\end{figure}

\begin{figure}
    \centering
     \caption{Historical Shock Decomposition for Private Consumtpion Growth Rate}
    \includegraphics [scale=.4]{Figures/gam_C_obs}
    \smallskip
    \captionsetup{width=.8\linewidth}
    \caption*{{\tiny The black line is the observed variable. The y-axis is percentage point deviation from steady state and the x-axis is quarters.  }}
\end{figure}

\begin{figure}
    \centering
     \caption{Historical Shock Decomposition for Total Investment Growth Rate}
    \includegraphics [scale=.4]{Figures/gam_I_obs}
    \smallskip
    \captionsetup{width=.8\linewidth}
    \caption*{{\tiny The black line is the observed variable. The y-axis is percentage point deviation from steady state and the x-axis is quarters.  }}
\end{figure}

\begin{figure}
    \centering
     \caption{Historical Shock Decomposition for Government Consumption Growth Rate}
    \includegraphics [scale=.4]{Figures/gam_G_obs}
    \smallskip
    \captionsetup{width=.8\linewidth}
    \caption*{{\tiny The black line is the observed variable. The y-axis is percentage point deviation from steady state and the x-axis is quarters.  }}
\end{figure}

\begin{figure}
    \centering
     \caption{Historical Shock Decomposition for Government Investment Growth Rate}
    \includegraphics [scale=.4]{Figures/gam_Ig_obs}
    \smallskip
    \captionsetup{width=.8\linewidth}
    \caption*{{\tiny The black line is the observed variable.The y-axis is percentage point deviation from steady state and the x-axis is quarters.  }}
\end{figure}

\begin{figure}
    \centering
     \caption{Historical Shock Decomposition for Transfers Growth Rate}
    \includegraphics [scale=.4]{Figures/gam_TR_obs}
    \smallskip
    \captionsetup{width=.8\linewidth}
    \caption*{{\tiny The black line is the observed variable. The y-axis is percentage point deviation from steady state and the x-axis is quarters.}}
\end{figure}

\begin{figure}
    \centering
     \caption{Historical Shock Decomposition for Hours}
    \includegraphics [scale=.4]{Figures/hours_obs}
    \smallskip
    \captionsetup{width=.8\linewidth}
    \caption*{{\tiny The black line is the observed variable. The y-axis is percent deviation from steady state and the x-axis is quarters.}}
\end{figure}

\begin{figure}
    \centering
     \caption{Historical Shock Decomposition for Oil Price}
    \includegraphics [scale=.4]{Figures/pCostar_obs}
    \smallskip
    \captionsetup{width=.8\linewidth}
    \caption*{{\tiny The black line is the observed variable.The y-axis is percent deviation from steady state and the x-axis is quarters.}}
\end{figure}

\begin{figure}
    \centering
     \caption{Historical Shock Decomposition for Trade Balance to GDP Ratio}
    \includegraphics [scale=.4]{Figures/ratio_TBY_obs}
    \smallskip
    \captionsetup{width=.8\linewidth}
    \caption*{{\tiny The black line is the observed variable. The y-axis is percentage point deviation from steady state and the x-axis is quarters.}}
\end{figure}

\begin{figure}
    \centering
     \caption{Historical Shock Decomposition for Foreign Interest Rate}
    \includegraphics [scale=.4]{Figures/Rstar_obs}
    \smallskip
    \captionsetup{width=.8\linewidth}
    \caption*{{\tiny The black line is the observed variable. The y-axis is percent deviation from steady state and the x-axis is quarters.}}
\end{figure}

\begin{figure}
    \centering
     \caption{Historical Shock Decomposition for Country Premium}
    \includegraphics [scale=.4]{Figures/xi_obs}
    \smallskip
    \captionsetup{width=.8\linewidth}
    \caption*{{\tiny The black line is the observed variable. The y-axis is percent deviation from steady state and the x-axis is quarters.}}
\end{figure}

\begin{figure}
    \centering
     \caption{Historical Shock Decomposition for Oil Quantity}
    \includegraphics [scale=.4]{Figures/YCo_obs}
    \smallskip
    \captionsetup{width=.8\linewidth}
    \caption*{{\tiny The black line is the observed variable. The y-axis is percent deviation from steady state and the x-axis is quarters.}}
\end{figure}

\begin{figure}
    \centering
     \caption{Historical Shock Decomposition for Foreign GDP} \label{fig:shockdecompforeigngdp}
    \includegraphics [scale=.4]{Figures/Ystar_obs}
    \smallskip
    \captionsetup{width=.8\linewidth}
    \caption*{{\tiny The black line is the observed variable. The y-axis is percent deviation from steady state and the x-axis is quarters.}}
\end{figure}

\section{Impulse-Response Analysis}\label{appendix:impulseresponses}

The next six subsections present simulations in which each of the six fiscal instruments is adjusted to achieve a reduction in the fiscal deficit-to-GDP ratio of 1 percentage point on impact. All simulations are performed under perfect foresight.

\subsection{An Increase in the Consumption Tax Rate}\label{app:increaseconsumptiontax}

Figure \ref{irf_tconsumption} presents a fiscal consolidation based on an increase in the consumption tax rate (plotted in the second row and first column of the figure; panel 2-1). The size of the tax rate increase, 1.6 pp above steady state, is designed to reduce the fiscal deficit-to-GDP ratio by 1 pp on impact (top-left panel).\footnote{We calibrate the consumption tax rate change using the government budget constraint in steady state. Explicitly, we solve the following equation for $x$:
$(pg^c + pg^i + tr + \xi R^*b^* - b^* - \chi p^{Co*}y^{Co*} - \tau^n w(1-\omega)h^R - \tau^k(r^k-\delta)\frac{k}{a} - (\tau^c+x)pc)/y = 
 (pg^c + pg^i + tr + \xi R^*b^* - b^* - \chi p^{Co*}y^{Co*} - \tau^n w(1-\omega)h^R - \tau^k(r^k-\delta)\frac{k}{a} - \tau^cpc)/y -.01 
\implies x = \frac{.01y}{pc}$. The deficit-to-GDP ratio does not fall by exactly 1 pp on impact because the calibration of the consumption tax rate assumed all other variables remained at steady state and thus did not account for the endogenous response of the economy to the tax change.} The simulation is a 5-year forecast, conditioning on the consumption tax rate lying 1.6 pp above its steady state and all other fiscal instruments remaining at their steady state for 20 quarters.

The decline in the fiscal deficit-to-GDP ratio generates a gradual decline in the (annualized) government debt-to-GDP ratio (panel 1-2). Aggregate private consumption declines in response to the higher consumption tax rate (panel 2-4). The decline in consumption is more immediate for the hand-to-mouth households (panel 4-2) and more gradual for the Ricardian households (panel 4-1), who can smooth consumption by increasing debt (panel 1-3). The increase in private debt is not sufficient to compensate the decline in public debt, so the aggregate foreign debt-to-GDP ratio declines (panel 1-4), generating downward pressure on the country risk premium (panel 2-2).

Ricardian households anticipate the decline in the country premium due to the fall in government debt which increases the present value of investment (higher Tobin's q), leading to higher private investment throughout the simulation horizon (panel 3-3). During the first couple of years, as the model economy adjusts to the higher consumption tax rate, non-oil GDP declines (panel 2-3). However, as private investment gradually expands in response to the lower country risk premium, output eventually increases above trend.

\begin{figure}[ht]
    \centering
     \caption{Consolidation: An Increase in the Consumption Tax Rate}
    \includegraphics [scale=.75]{Consumption_Tax_IRF.eps}
    \label{irf_tconsumption}
    \smallskip
    \captionsetup{width=.6\linewidth}
    \caption*{\footnotesize{\emph{Note}: Fiscal consolidation designed to reduce the fiscal deficit-to-GDP ratio by approximately 1 percentage point on impact through an increase in the consumption tax rate, holding all other fiscal instruments constant. X axes denote quarters. Y axes denote percent deviations from steady state, except for variables expressed as ratios to GDP and the country premium, which are expressed as percentage point deviation from steady state. Public and private debt, and the country premium, are annualized.}}
    \end{figure}
    
\subsection{An Increase in the Labor Income Tax Rate}\label{app:increaselabortax}

Figure \ref{irf_tlabor} presents a fiscal consolidation based on an increase in the labor income tax rate (plotted in the second row and first column; panel 2-1). Again, the increase in the tax rate, of 5.6 pp above steady state, is calibrated to generate a decline in the fiscal deficit-to-GDP ratio of approximately 1 pp on impact (top-left panel).\footnote{We calibrate the labor tax rate change using the government budget constraint in steady state. Explicitly, we solve the following equation for $x$:
$(pg^c + pg^i + tr + \xi R^*b^* - b^* - \chi p^{Co*}y^{Co*} - (\tau^n+x) w(1-\omega)h^R - \tau^k(r^k-\delta)\frac{k}{a} - \tau^cpc)/y = 
 (pg^c + pg^i + tr + \xi R^*b^* - b^* - \chi p^{Co*}y^{Co*} - \tau^n w(1-\omega)h^R - \tau^k(r^k-\delta)\frac{k}{a} - \tau^cpc)/y -.01 \implies x = \frac{.01y}{w(1-\omega)h^R}$.} The fall in the fiscal deficit leads to a gradual decline in the (annualized) government debt-to-GDP ratio (panel 1-2).

The higher labor income tax rate generates a decline in hours worked (panel 3-2) and an increase in the real wage due to its contractionary effect on labor supply (panel 3-1). As the labor income tax is only levied on Ricardian households, the decline in aggregate hours is driven by this group (panel 4-3). These households also reduce consumption gradually but substantially in response to the intratemporal substitution effect induced by lower labor supply (panel 4-1). The group of hand-to-mouth households, on the other hand, increases consumption in response to the rise in their disposable income due to the higher wage (panel 4-2).

Aggregate private consumption declines (panel 2-4), driven by the behavior of Ricardian households, despite them issuing private foreign debt (panel 1-3) in order to reduce consumption gradually. The aggregate foreign debt-to-GDP ratio (panel 1-4), however, declines, driven by public debt, so the country risk premium also falls (panel 2-2). The lower premium induces an increase in private investment (panel 3-3). All in all, non-oil GDP (panel 2-3) falls below trend, driven by the decline in hours worked despite the higher capital stock.
    
\begin{figure}[ht]
    \centering
     \caption{Consolidation: An Increase in the Labor Income Tax Rate}%IRF to Labor Income Tax}
    \includegraphics [scale=.75]{Figures/Labor_Tax_IRF.eps}
    \label{irf_tlabor}
    \smallskip
    \captionsetup{width=.6\linewidth}
    \caption*{\footnotesize{\emph{Note}: Fiscal consolidation designed to reduce the fiscal deficit-to-GDP ratio by approximately 1 percentage point on impact through an increase in the labor income tax rate, holding all other fiscal instruments constant. For more details, see the note to figure \ref{irf_tconsumption}.}}
\end{figure}

\subsection{An Increase in the Capital Income Tax Rate}\label{app:increasecapitaltax}

Figure \ref{irf_tcapital} presents a fiscal consolidation based on an increase in the capital income tax rate (panel 2-1) designed to reduce the fiscal deficit-to-GDP ratio by approximately 1 pp on impact (top-left panel). To achieve this, the capital income tax rate, which is levied only on the group of Ricardian households, has to increase by 6.5 pp above steady state.\footnote{We calibrate the capital tax rate change using the government budget constraint in steady state. Explicitly, we solve the following equation for $x$:
$(pg^c + pg^i + tr + \xi R^*b^* - b^* - \chi p^{Co*}y^{Co*} - \tau^n w(1-\omega)h^R - (\tau^k+x)(r^k-\delta)\frac{k}{a} - \tau^cpc)/y = 
 (pg^c + pg^i + tr + \xi R^*b^* - b^* - \chi p^{Co*}y^{Co*} - \tau^n w(1-\omega)h^R - \tau^k(r^k-\delta)\frac{k}{a} - \tau^cpc)/y -.01 
\implies x = \frac{.01y}{(r^k-\delta)\frac{k}{a}}$.} The decline in the fiscal deficit leads to a gradual decline in the (annualized) government debt-to-GDP ratio (panel 1-2).

The country risk premium falls in response to lower aggregate debt (panel 2-2) which should lead to an increase in investment projects. However, private investment declines substantially (panel 3-3) because the force of the higher capital income tax rate dominates the decline in the country premium. The lower risk premium generates an intertemporal substitution effect that leads Ricardian households to increase consumption (panel 4-1), which drives a slight increase in aggregate private consumption (panel 2-4). The decline in investment affects the capital stock, which lowers the marginal product of labor and thus the real wage (panel 3-1), lowering consumption by hand-to-mouth households. Overall, the higher capital income tax generates a persistent decline in non-oil GDP (panel 2-3).

\begin{figure}[ht]
    \centering
     \caption{Consolidation: An Increase in the Capital Income Tax Rate}%IRF to Capital Income Tax}
    \includegraphics [scale=.75]{Figures/Capital_Tax_IRF.eps}
    \label{irf_tcapital}
    \smallskip
    \captionsetup{width=.6\linewidth}
    \caption*{\footnotesize{\emph{Note}: Fiscal consolidation designed to reduce the fiscal deficit-to-GDP ratio by approximately 1 percentage point on impact through an increase in the capital income tax rate, holding all other fiscal instruments constant. For more details, see the note to figure \ref{irf_tconsumption}.
    }}
\end{figure}

\subsection{A Reduction in Government Consumption}\label{app:reductionconsumption}

We now turn to fiscal consolidation scenarios based on adjustments of the components of government expenditure. Figure \ref{irf_consumption} presents a fiscal consolidation based on a reduction in government consumption. Again, we calibrate the size of the adjustment so that it generates a 1 pp decline, on impact, in the fiscal deficit-to-GDP ratio (top-left panel in the figure). We find that a 9\% reduction in the level of government consumption generates the required adjustment (panel 2-1).\footnote{We calibrate the government consumption change using the government budget constraint in steady state. Explicitly, we solve the following equation for $x$:
$(p(g^c+x) + pg^i + tr + \xi R^*b^* - b^* - \chi p^{Co*}y^{Co*} - \tau^n w(1-\omega)h^R - \tau^k(r^k-\delta)\frac{k}{a} - \tau^cpc)/y = 
 (pg^c + pg^i + tr + \xi R^*b^* - b^* - \chi p^{Co*}y^{Co*} - \tau^n w(1-\omega)h^R - \tau^k(r^k-\delta)\frac{k}{a} - \tau^cpc)/y -.01 
\implies x =  \frac{.01y}{p}$.}

The contraction in government consumption gradually reduces the government debt-to-GDP ratio (panel 1-2), which in turn leads to a substantial decline in the country risk premium (more than 30 basis points after five years; panel 2-2). The intertemporal substitution effect induced by the lower country premium leads the Ricardian households to increase consumption (panel 4-1). The decline in the country premium also induces an increase in investment (panel 3-3) due to the increase in its present value. The group of hand-to-mouth households reduces consumption initially (panel 4-2), because of the complementarity between public and private consumption and their inability to smooth the adjustment by issuing foreign debt, which is what the Ricardian households do (panel 1-3). Eventually, however, the higher stock of private capital pushes the marginal product of labor and the real wage upwards (panel 3-1), which allows the hand-to-mouth households to increase consumption. Hours worked by both groups decline (panels 4-3 and 4-4). This is another consequence of the complementarity between public and private consumption\textemdash the decline in government consumption puts downward pressure on the marginal utility of private consumption, which leads to an intratemporal substitution towards leisure. Overall, the lower aggregate hours worked (panel 3-2) generate an initial contraction in non-oil GDP (panel 2-3), but the expansion in the capital stock eventually pulls it above trend.

\begin{figure}[ht]
    \centering
     \caption{Consolidation: A Reduction in Government Consumption}
    \includegraphics [scale=.8]{Figures/Government_Consumption_IRF.eps}
    \label{irf_consumption}
    \smallskip
    \captionsetup{width=.6\linewidth}
    \caption*{\footnotesize{\emph{Note}: Fiscal consolidation designed to reduce the fiscal deficit-to-GDP ratio by approximately 1 percentage point on impact through a reduction in government consumption, holding all other fiscal instruments constant. For more details, see the note to figure \ref{irf_tconsumption}.}}
\end{figure}

\subsection{A Reduction in Government Investment}\label{app:reductioninvestment}

Figure \ref{irf_investment} presents a consolidation scenario based on a reduction in government investment. To generate a 1 pp decline in the fiscal deficit-to-GDP ratio (top-left panel), the level of government investment has to decline 37\% (panel 2-1).\footnote{We calibrate the government investment change using the government budget constraint in steady state. Explicitly, we solve the following equation for $x$:
$(pg^c + p(g^i+x) + tr + \xi R^*b^* - b^* - \chi p^{Co*}y^{Co*} - \tau^n w(1-\omega)h^R - \tau^k(r^k-\delta)\frac{k}{a} - \tau^cpc)/y = 
 (pg^c + pg^i + tr + \xi R^*b^* - b^* - \chi p^{Co*}y^{Co*} - \tau^n w(1-\omega)h^R - \tau^k(r^k-\delta)\frac{k}{a} - \tau^cpc)/y -.01 
\implies x  = \frac{.01y}{p}$.}

The reduction in government investment leads to a substantial decline in the country risk premium (panel 2-2) through its effect on the government debt-to-GDP ratio (panel 1-2), which generates an intertemporal substitution effect for the Ricardian households. These households increase consumption (panel 4-1) and leisure (panel 4-3). The decline in the premium also induces an increase in private investment (panel 3-3). However, the decline in the stock of public capital is larger so the marginal product of labor and the real wage decline persistently (panel 3-1). Due to the reduction in the real wage, the hand-to-mouth households have no alternative but to reduce consumption. 

The decline in government capital and aggregate hours worked (panel 3-2) weigh on non-oil GDP which remains below trend throughout the five-year horizon (panel 2-3), despite the higher private capital stock.

\begin{figure}[ht]
    \centering
     \caption{Consolidation: A Reduction in Government Investment}%IRF to Government Investment}
    \includegraphics [scale=.75]{Government_Investment_IRF.eps}
    \label{irf_investment}
    \smallskip
    \captionsetup{width=.6\linewidth}
    \caption*{\footnotesize{\emph{Note}: Fiscal consolidation designed to reduce the fiscal deficit-to-GDP ratio by approximately 1 percentage point on impact through a reduction in government investment, holding all other fiscal instruments constant. For more details, see the note to figure \ref{irf_tconsumption}.
    }}
\end{figure}

\subsection{A Reduction in Government Transfers}\label{app:reductiontransfers}

We now discuss a fiscal consolidation scenario based on a reduction of the final component of government expenditure, namely transfers, which we assume are only paid to the hand-to-mouth households. The results are shown in figure \ref{irf_transfers}. To generate a 1-pp reduction in the fiscal deficit-to-GDP ratio on impact (panel 1-1), the level of transfers must decline 20\%, which is equivalent to 1 pp of GDP (panel 2-1).\footnote{We calibrate the transfer change using the government budget constraint in steady state. Explicitly, we solve the following equation for $x$:
$(pg^c + pg^i + (tr+x_ + \xi R^*b^* - b^* - \chi p^{Co*}y^{Co*} - \tau^n w(1-\omega)h^R - \tau^k(r^k-\delta)\frac{k}{a} - \tau^cpc)/y = 
 (pg^c + pg^i + tr + \xi R^*b^* - b^* - \chi p^{Co*}y^{Co*} - \tau^n w(1-\omega)h^R - \tau^k(r^k-\delta)\frac{k}{a} - \tau^cpc)/y -.01 
\implies x = .01y$.} 

The reduction in transfers is a substantial income shock to the hand-to-mouth households, who reduce consumption and leisure drastically in response (panels 4-2 and 4-4). The improvement in the fiscal deficit gradually lowers the government debt-to-GDP ratio (panel 1-2), which generates a substantial decline in the country risk premium (panel 2-2). The Ricardian households, who are not directly affected by the reduction in transfers, increase consumption and leisure because of the intertemporal substitution effect caused by the lower interest rate (panels 4-1 and 4-3). The decline in the country premium also generates an increase in private investment (panel 3-3). The increase in aggregate labor supply (panel 3-2) puts downward pressure on the real wage (panel 3-1). The overall expansion in hours worked and the capital stock induce an expansion in non-oil GDP (panel 2-3) above trend.

\begin{figure}[ht]
    \centering
     \caption{Consolidation: A Reduction in Government Transfers}%IRF to Transfers}
    \includegraphics [scale=.75]{Figures/Transfers_IRF}
    \label{irf_transfers}
    \smallskip
    \captionsetup{width=.6\linewidth}
    \caption*{\footnotesize{\emph{Note}: Fiscal consolidation designed to reduce the fiscal deficit-to-GDP ratio by approximately 1 percentage point on impact through a reduction in government transfers, holding all other fiscal instruments constant. For more details, see the note to figure \ref{irf_tconsumption}.
    }}
\end{figure}

\section{The Effect of Capital Controls}\label{app:capitalcontrols}
We evaluate the effect of the fiscal consolidation in a country in which the pass-through from external rates and the country premium is only partial, so that the benefits of lowering borrowing costs from reducing government debt are not as envisioned in the setup of our model. This feature could be common to certain commodity-exporting countries that have capital control measures in place, such as Ecuador. To that end, we assume that the return on foreign bonds that private agents receive or the interest rate they pay on borrowing is not $r^*_t = R^*_{t-1}\xi_{t-1}$, as specified in \eqref{eq:interestrate}. To allow for only partial pass-through from either $R^*_{t-1}$ or $\xi_{t-1}$, we specify the relevant interest rate that private agents face for consumption and investment decisions as $\tilde{r}^*_t = R^*\xi + \mu (R^*_{t-1}\xi_{t-1}-R^*\xi)$, with $\mu=0.15$.\footnote{Notice that $\mu=1$ implies complete pass-through and the absence of capital controls.}
\begin{figure}[ht]
    \centering
     \caption{Full Fiscal Consolidation with Capital Controls}
    \includegraphics [scale=.8]{Figures/IMF_Simulation_capital_control.eps}
    \label{fig:capitalcontrols}
    \smallskip
    \captionsetup{width=.8\linewidth}
    \caption*{\footnotesize{\emph{Note}: Fiscal consolidation combining adjustments of multiple revenue and expenditure instruments designed to approximate the agreement between the government of Ecuador and the IMF for the 2020\textendash 2025 period, as shown in tables \ref{tIMF_taxes} and \ref{tIMF_spending}. X axes denote quarters. Y axes denote  percent deviations from steady state, except for variables expressed as ratios to GDP and the country premium, which are expressed as percentage point deviation from steady state. The country premium is annualized. We assume that that the interest rate that face private domestic agents is $\tilde{r}^*_t = R^*\xi + \mu (R^*_{t-1}\xi_{t-1}-R^*\xi)$, with $\mu=0.15$. The continuous blue lines depict the effects of the fiscal consolidation absent capital controls while the dashed orange lines illustrate the effects under capital controls.}}
\end{figure}

\section{Accounting for Parameter Uncertainty}\label{app:uncertainty}

This appendix assesses the sensitivity of our fiscal-consolidation results to parameter uncertainty. We report transition paths under the baseline oil-revenue trajectory together with Bayesian credible intervals constructed from posterior draws of the structural parameters. 

Let $\theta$ denote the vector of structural parameters. We take the posterior distribution $p(\theta \mid \mathcal{D})$ obtained from the linearized stochastic estimation (Section~\ref{app:estimatedparameters}) and consider deterministic transition dynamics under perfect foresight as in Section~\ref{sec:fiscalpolicy}. For each draw $\theta^{(j)}$, we solve the nonlinear model with: (i) the IMF-supported program’s quarter-by-quarter fiscal-instrument paths imposed as deterministic sequences over the 24-quarter consolidation window, (ii) the baseline oil-revenue trajectory, and (iii) the same initial conditions and terminal assumptions as in the baseline simulation (see Section~\ref{sec:fiscalpolicy}). We draw $J=200$ parameter vectors $\{\theta^{(j)}\}_{j=1}^{J}$ from $p(\theta \mid \mathcal{D})$ (thinning the MCMC output to reduce autocorrelation). For each draw, we solve the perfect-foresight transition using \texttt{perfect\_foresight\_setup} and \texttt{perfect\_foresight\_solver} in Dynare. We then construct pointwise (across draws) $90\%$ credible intervals for each variable $x_t$ by taking the 5th and 95th percentiles of $\{x_t^{(j)}\}_{j=1}^{J}$ for every quarter $t=1,\dots,24$. The solid line in the figures corresponds to the posterior-mean parameter vector.

The credible intervals indicate economically meaningful uncertainty around the magnitude and timing of the responses, but the qualitative dynamics of the fiscal-consolidation experiment are robust across posterior draws. 

In particular: (a) non-oil output exhibits a modest trough during the early program years and recovers thereafter; and (b) the deficit-to-GDP ratio improves persistently. The posterior bands remain on the same side of zero as the posterior-mean paths at relevant horizons for most variables, confirming that the directional results emphasized in the main text are not driven by a particular parametrization.

\begin{figure}[ht]
    \centering
     \caption{Baseline Fiscal Consolidation with 90\% Credible Interval}
    \includegraphics [scale=1.2]{Figures/IMF_Simulation_bayesian.eps}
    \label{fig:irf_bayesian}
    \smallskip
    \captionsetup{width=0.9\linewidth}
    \caption*{\footnotesize{\emph{Note}: Solid line: transition path at posterior-mean parameters. Dashed lines: pointwise $90\%$ Bayesian credible interval across $J=200$ posterior draws. X axes denote quarters. Y axes denote  percent deviations from steady state, except for variables expressed as ratios to GDP and the country premium, which are expressed as percentage point deviation from steady state. Instrument paths and oil revenue follow the program and hence show no posterior dispersion during the consolidation window.}}
\end{figure}


\end{document}